\chapter{\texorpdfstring{The $\beta\gamma$ system}{The beta-gamma system}}\label{chap:beta-gamma}
% Regime I: quadratic (Convention~\ref{conv:regime-tags}).

A chiral algebra with a weight-zero generator violates every
positive-grading hypothesis standard in the theory of vertex
operator algebras. At $\lambda = 0$ the $\beta\gamma$ system has a
generator $\gamma$ of conformal weight zero: the mode expansion
$\gamma(z) = \sum_n \gamma_n z^{-n}$ contains a zero mode
$\gamma_0$ that acts on the vacuum without raising energy, the
character diverges as $1/(1-q)$ rather than as a $q$-series with
positive coefficients, and standard positive-energy filtrations
collapse. The algebra cannot be treated as a Heisenberg-type or
Kac--Moody-type object, where a positive grading forces the shadow
tower to terminate at degree~$2$ or~$3$.

The depth classification assigns the standard conformal-weight family
$\beta\gamma_\lambda$ to class~$\mathsf{C}$. Heisenberg and lattice
algebras (class~$\mathsf{G}$) have shadow towers terminating at
degree~$2$; affine Kac--Moody algebras (class~$\mathsf{L}$) terminate
at degree~$3$; Virasoro and $\mathcal{W}_N$ (class~$\mathsf{M}$) run
to infinity. Class~$\mathsf{C}$ is the contact/quartic class: in the
standard landscape it is represented by the Koszul pair
$\beta\gamma_\lambda \leftrightarrow bc_\lambda$, not by a single
one-dimensional primary line. The standard-family coefficients are
\[
S_2 = 6\lambda^2 - 6\lambda + 1,\qquad
S_3 = 0,\qquad
S_4 = -\tfrac{5}{12},\qquad
S_r = 0\quad (r\ge 5),
\]
so the shadow depth is $r_{\max}=4$. By contrast, the internal
weight-changing line is shadow-trivial; its scalar quartic invariant
$\mu_{\beta\gamma}$ vanishes by rank-one abelian rigidity.

Three invariants separate the local OPE from the shadow tower:
the maximal generator OPE pole order is $p_{\max}=1$, the
collision-residue depth is $k_{\max}=p_{\max}-1=0$, and the
standard-family shadow depth is $r_{\max}=4$. The collision
residue therefore vanishes even though the standard family carries
class~$\mathsf{C}$ contact data. The Koszul dual algebra is the
$bc$ ghost system (Theorem~\ref{thm:betagamma-fermion-koszul}),
exchanging bosonic and fermionic statistics: the chiral lift of
$\mathrm{Sym}(V)^! = \Lambda(V^*)$. The central charges cancel,
$c_{\beta\gamma}+c_{bc}=0$, and hence
$\kappa(\beta\gamma_\lambda)+\kappa(bc_\lambda)=0$ on the scalar
surface. Two generators of complementary conformal weight $\lambda$ and
$1-\lambda$, coupled by the simple-pole OPE
$\beta(z)\gamma(w)\sim 1/(z{-}w)$, are the minimal data.

\begin{table}[ht]
\centering
\small
\caption{Five-theorem verification for
$\beta\gamma$ at generic conformal weight
$\lambda$.}\label{tab:betagamma-five-theorems}
\begin{tabular}{lll}
\toprule
\textbf{Theorem} & \textbf{Statement for $\beta\gamma$}
 & \textbf{Anchor} \\
\midrule
A (adjunction) &
 $A^{\mathrm i}=H^\bullet(\bar B(A))$ and
 $A^!\simeq bc_\lambda$ for $A=\beta\gamma_\lambda$
 & Thm~\ref{thm:betagamma-bc-koszul} \\
B (inversion) &
 $\Omega(\barB(\beta\gamma)) \xrightarrow{\sim}
 \beta\gamma$
 & Prop~\ref{prop:betagamma-bar-acyclicity} \\
C (complementarity) &
 $\kappa(\beta\gamma) + \kappa(bc) = 0$
 & \S\ref{sec:koszul-duality-free-fields} \\
D (modular char.) &
 $\kappa(\beta\gamma)
 = c/2 = 6\lambda^2{-}6\lambda{+}1$
 & Rem~\ref{rem:betagamma-shadow-archetype-free} \\
H (chiral Hochschild) &
 $\ChirHoch^n(\beta\gamma_\lambda)=0$ for
 $n\notin\{0,1,2\}$
 & Thm~\ref{thm:hochschild-polynomial-growth} \\
\bottomrule
\end{tabular}
\end{table}

\begin{table}[ht]
\centering
\small
\caption{Shadow archetype data for
$\beta\gamma$.}\label{tab:betagamma-shadow-archetype}
\begin{tabular}{ll}
\toprule
\textbf{Invariant} & \textbf{Value} \\
\midrule
Class & C (contact/quartic) \\
Shadow depth $r_{\max}$ & $4$ \\
$\kappa(\beta\gamma)$ &
 $c/2 = 6\lambda^2 - 6\lambda + 1$ \\
Cubic shadow $\mathfrak{C}$ &
 $S_3=0$ on the standard family; mixed OPE gives the bar differential \\
Quartic $Q^{\mathrm{contact}}$ &
 $S_4=-5/12$ on the standard family; $\mu_{\beta\gamma}=0$ on the weight line \\
Quintic and higher & $S_r=0$ for $r\ge5$ on the standard family \\
$r$-matrix &
 $r^{\mathrm{coll}}_{\beta\gamma}(z)=r^{\mathrm{coll}}_{bc}(z)=0$;
 ordered contact and line transports are separate \\
Koszul dual & $bc$ \\
Complementarity & $\kappa(\beta\gamma) + \kappa(bc) = 0$ \\
$L_\infty$-formality &
 standard-family contact datum; weight line formally quadratic \\
\bottomrule
\end{tabular}
\end{table}

\noindent
The five objects remain distinct throughout this chapter. For
$A=\beta\gamma_\lambda$, $B(A)$ is the ordered bar coalgebra,
$A^{\mathrm{i}}=H^\bullet(B(A))$ is the Koszul-dual coalgebra,
$A^!=(A^{\mathrm{i}})^\vee\simeq bc_\lambda$ is the Verdier/Koszul
dual algebra on finite-type truncations, and
$Z^{\mathrm{der}}_{\mathrm{ch}}(A)=\ChirHoch^\bullet(A,A)$ is the
cochain-level chiral closed sector. The identity $\Omega(B(A))\simeq A$ is
bar--cobar inversion; it is not the assertion $B(A)=A^!$, and it is
not a derived-centre computation. Theorem~H is the concentration of
this last object in degrees $\{0,1,2\}$, not a statement about the
bare bar coalgebra or about the Koszul-dual algebra.

\noindent
The modular Koszul triple
(Definition~\ref{def:modular-koszul-triple}) is
\begin{equation}\label{eq:betagamma-triple-ch}
\mathfrak{T}_{\beta\gamma}
\;=\;
\bigl(\,\beta\gamma_\lambda,\;\;
bc_\lambda,\;\;
r^{\mathrm{coll}}(z) = 0\,\bigr).
\end{equation}
The Koszul duality $(\beta\gamma)^! \simeq bc$ exchanges
boson--fermion statistics
(Theorem~\ref{thm:betagamma-bc-koszul}). The collision-residue
$r$-matrix vanishes: the $\beta\gamma$ OPE
$\beta(z)\gamma(w) \sim 1/(z{-}w)$ has only a simple pole,
which is absorbed by $d\log$ extraction. The same convention gives
$r^{\mathrm{coll}}_{bc}(z)=0$ for the $bc$ ghost system. The leading
standard-family interaction datum is the quartic contact coefficient
$S_4=-5/12$, not the collision $r$-matrix. If the ordered bar
transfer retains the regular contact operators, the tensors are
\[
\Theta_{\beta\gamma}^{\mathrm{ord}}
=\beta\otimes\gamma-\gamma\otimes\beta,
\qquad
\Theta_{bc}^{\mathrm{ord}}
=b\otimes c+c\otimes b .
\]
The second sign is the graded Clifford sign; after desuspension to an
ordinary tensor algebra the same class may be written with a minus.
The graded formula above is the convention used here. In either
statistical sector $\Theta^2=0$ and
$R_{\mathrm{contact}}=\id+2\pi i\,\Theta$; this unipotent transport
is not a pole-valued collision residue.

The genus-$0$ curvature vanishes ($m_0 = 0$: the $\beta\gamma$ OPE
has no double pole), so the genus-$0$ bar complex is uncurved.
On the weight-changing line, the degree-$2$ shadow vanishes
($\Theta_{\beta\gamma}^{\leq 2}\big|_{\mathrm{w.c.}} = 0$),
though the global modular characteristic
$\kappa(\beta\gamma) = c/2$ is generically nonzero.
The standard conformal-weight family has nonzero quartic contact
coefficient $S_4=-5/12$ and no tail beyond degree~$4$. The
weight-changing line is the rank-one degeneration of this picture:
there $\Theta_{\beta\gamma}^{\le r}|_{\mathrm{w.c.}}=0$ for
$r\ge2$ and $\mu_{\beta\gamma}=0$. Thus class~C
($r_{\max}=4$) belongs to the standard family, not to the
one-dimensional weight line.

\begin{remark}[Shadow obstruction tower for
$\beta\gamma$]\label{rem:betagamma-master-mc}
For the $\beta\gamma$ system, $m_0 = 0$ (the bar complex is
uncurved, since the $\beta\gamma$ OPE has no double pole) but the shadow obstruction tower is
nontrivial beyond the scalar level.
On the weight-changing line (the $\lambda$-marginal slice),
$\Theta_{\beta\gamma}^{\leq 2}\big|_{\mathrm{w.c.}} = 0$
and in fact
$\Theta_{\beta\gamma}^{\leq r}\big|_{\mathrm{w.c.}} = 0$ for
$r\ge2$. The quartic contact datum of class~C belongs to the full
standard conformal-weight family, where $S_4=-5/12$; its scalar
projection on the rank-one weight-changing line is
$\mu_{\beta\gamma}=0$ by rank-one abelian rigidity.
The global modular characteristic
$\kappa(\beta\gamma) = c/2 = 6\lambda^2 - 6\lambda + 1$
is generically nonzero
(Proposition~\ref{prop:betagamma-obstruction-coefficient});
the degree-$2$ vanishing is specific to the weight-changing slice,
not a statement about $\kappa$.
The distinction between the standard family and the rank-one line is
the reason that $r_{\max}=4$ and $\mu_{\beta\gamma}=0$ are compatible.
\end{remark}

\begin{remark}[Three-pillar interpretation: $\beta\gamma$]
\label{rem:betagamma-three-pillar}
\index{beta-gamma@$\beta\gamma$!three-pillar interpretation}
In the three-pillar architecture
(\S\ref{sec:concordance-three-pillars}):
(i)~the mixed $\beta\gamma$ OPE supplies the bar differential, while
the standard-family shadow has $S_3=0$ and $S_4=-5/12$.
The rank-one weight-changing scalar
$\mu_{\beta\gamma}=0$ by rank-one abelian rigidity
(Corollary~\ref{cor:nms-betagamma-mu-vanishing});
(ii)~the convolution $sL_\infty$-algebra separates the full
standard-family contact datum from its formally quadratic
rank-one weight-line restriction. The shadow algebra
$(\beta\gamma)^{\mathrm{sh}}$ is class~C for the standard family
and trivial on the internal weight-changing line;
(iii)~the log-FM geometry $\overline{\operatorname{FM}}_n(C|D)$ is
relevant when punctures from the $\gamma$-insertion are present:
the quartic contact class arises as a codimension-$2$ residue on the
log-FM boundary where four insertions collide near a puncture.
\end{remark}

\begin{remark}[Genus-two shell profile: $\beta\gamma$]
\label{rem:betagamma-genus-two-shells}
\index{beta-gamma@$\beta\gamma$!genus-two shells}
\index{genus-two shells!contact/quartic archetype}
In the genus-two shell decomposition
(Construction~\ref{const:vol1-genus-two-shells}),
the rank-one weight-changing branch sees only the first shell:
$\Theta^{(2)}_{\mathrm{loop} \circ \mathrm{loop}}$. The
separating/loop shell vanishes there because strict formality at tree
level kills iterated cubic-loop interactions, and the planted-forest
shell has zero scalar projection because
$\mu_{\beta\gamma}=0$
(Corollary~\ref{cor:nms-betagamma-mu-vanishing}). On the standard
class-$\mathsf C$ surface the planted-forest contact channel is present:
it is controlled by $S_4=-5/12$, while the scalar tower still
terminates at degree~$4$.
\end{remark}

\begin{remark}[Four-level structure: $\beta\gamma$]
\label{rem:betagamma-four-level}
The $\beta\gamma$ system separates four levels of the modular
engine:
(i)~symplectic bosons on a curve;
(ii)~the strict local vertex-algebra OPE, with no higher local
$m_k$ for $k\ge3$;
(iii)~the transferred modular shadow calculus, where the standard
family has $S_3=0$ and $S_4=-5/12$;
(iv)~the rank-one weight-changing branch, where
$\mu_{\beta\gamma}=0$
(Corollary~\ref{cor:nms-betagamma-mu-vanishing}) and the branch
potential is quadratic.
\end{remark}

\begin{remark}[Chart-class enumeration: $\beta\gamma$]
\label{rem:betagamma-chart-class-enumeration}
\index{chart-class!beta-gamma@$\beta\gamma$}
\emph{Type signature:} (Open quadrant; chart presentation;
Beilinson level~$1$; HP = factorisation dg-cat
$\mathcal{C}^{\mathrm{op}}$ on $(X, D, \tau)$, vacuum $b$,
$A_b = \mathrm{End}_\mathcal{C}(b)$; archetype label invariant
under Morita equivalence on the level-$1\to 0$ fibre).
The chart algebra
\[
A_b^{\beta\gamma, \lambda} =
\mathrm{End}_{\mathcal{C}^{\mathrm{op}}_X}(b^{\beta\gamma,\lambda})
\;\simeq\;
\beta\gamma_\lambda\otimes_{\mathcal{D}_X}(\,\cdot\,)
\]
on the standard vacuum $b^{\beta\gamma,\lambda}$ realises the
two-generator first-order chiral algebra at conformal weight
$\lambda$. The five-archetype label is~$\mathsf{C}$ (contact/quartic,
$r_{\max}=4$). Three Morita-distinct chart presentations exist:
\begin{enumerate}[label=\textup{(C\arabic*)}]
\item \emph{Standard $\beta\gamma_\lambda$ chart}: generators
$\beta(z),\gamma(z)$ at conformal weights $\lambda, 1-\lambda$ with
simple-pole OPE $\beta(z)\gamma(w)\sim 1/(z{-}w)$. The standard
landscape census places this chart at the $\Phi_{\mathrm{ell}}$
locus (see Chapter~\ref{ch:landscape-census}).
\item \emph{Symplectic boson chart}: $\lambda=1/2$, the
weight-symmetric point with $\kappa(\beta\gamma_{1/2})=-1/2$.
Morita-equivalent to (C1) within the standard family by the
$\lambda\mapsto\lambda$ specialisation.
\item \emph{$bc$ Verdier-dual chart}: the Koszul-dual chart
$bc_\lambda$ exchanges Bose and Fermi statistics
(Theorem~\ref{thm:betagamma-bc-koszul}) but preserves shadow depth.
$bc_\lambda$ is a class-$\mathsf{C}$ chart in the same row of the
five-archetype matrix; its five-$\kappa$ tuple has the opposite
sign on the categorical and chiral-Hodge entries.
\end{enumerate}
The class-$\mathsf{C}$ label is chart-independent across (C1)--(C3)
on the standard conformal-weight family. The internal
weight-changing slice (where $\Theta_{\beta\gamma}^{\le r}=0$ for
$r\ge2$) is shadow-trivial and does not constitute an independent
chart class. Within Vol~III the matching $\Phi_d^{\mathrm{FA}}$
specialisation at $d=3$ on the K3$\times$E target locus
(level~$1\leftrightarrow 1$) yields the curved
$\beta\gamma$-system input to the two-stage CY-to-chiral functor.
\end{remark}

\begin{remark}[Theorem A in four ambients: $\beta\gamma$]
\label{rem:betagamma-theorem-A-four-lanes}
\index{Theorem A!four ambients!beta-gamma@$\beta\gamma$}
On the chart algebra $A_b^{\beta\gamma,\lambda}=\beta\gamma_\lambda$,
Theorem~A (Theorem~\ref{thm:A-infinity-2} in
Chapter~\ref{chap:theorem-A-infinity-2}) appears in the four ambients
$\mathrm{L}1,\mathrm{L}2,\mathrm{L}3,\mathrm{L}4$:
\begin{itemize}
\item \emph{$\mathrm{L}1$ (chain level, $\mathrm{Ch}(\mathrm{Vect})$).}
The bar--cobar adjunction is realised by the explicit
$\beta\gamma$ bar complex
$\bar{B}^{\mathrm{ch}}(\beta\gamma_\lambda)$ of
Theorem~\ref{thm:betagamma-bar-complex}; the witness chain
homotopy is the contact one $h_C=h_{\mathrm{LV}}/\mathcal N(\beta\gamma_\lambda)$,
where the level-$2$ Verdier-Ran-on-bar norm
$\mathcal N(\beta\gamma_\lambda)=1$ is the free-field OPE
normalisation (Convention~\ref{conv:A-two-kappa-shriek}), so
$h_C=h_{\mathrm{LV}}$. This denominator is the level-$2$ norm, not
the algebra-level complementarity ceiling
$K^\kappa=\kappa+\kappa^!=0$ of the $\beta\gamma$/$bc$ free-field
pair: the ceiling vanishes by Verdier sign-flip, while the homotopy
denominator $\mathcal N=1$ is finite.
\item \emph{$\mathrm{L}2$ (Quillen-equivalence).}
The Vallette bar--cobar Quillen equivalence
\cite[Theorem~2.1]{Val16} applies in the Francis--Gaitsgory
factorisation ambient on the entire conformal-weight family
$\lambda\in\mathbb{C}$ (no resonance); the Hinich~\cite{Hinich2003}
promotion is unobstructed.
\item \emph{$\mathrm{L}3$ ($(\infty,1)$-categorical).}
The $(\infty,1)$-localisation gives the adjoint equivalence between
$\Fact^{\aug}(X)$-objects in the class-$\mathsf{C}$ stratum and
their $\CoFact^{\conil,\comp}$-counterparts; the Lurie HA \S5.5
parallel topological case applies after the chiral-to-topological
specialisation.
\item \emph{$\mathrm{L}4$ ($(\infty,2)$-properad).}
Theorem~$A^{\infty,2}$ (Theorem~\ref{thm:A-infinity-2})
inscribes the $(\infty,2)$-categorical adjoint equivalence on the
properadic envelope generated by $\beta\gamma_\lambda$; the
Hackney--Robertson lift carries the quartic contact datum
$S_4=-5/12$ as the only non-Gaussian properadic correction.
\end{itemize}
The four ambients coincide on the entire $\lambda$-family because
$\beta\gamma_\lambda$ is Koszul for all $\lambda$
(no critical resonance analogous to the affine Kac--Moody
critical level).
\end{remark}

\begin{remark}[Five-$\kappa$ stratification: $\beta\gamma$]
\label{rem:betagamma-five-kappa-stratification}
\index{kappa stratification@$\kappa$-stratification!beta-gamma row C}
The five $\kappa$-measurements
$\{\kappa_{\mathrm{cat}}, \kappa^{\mathrm{Hodge}}_{\mathrm{ch}},
\kappa^{\mathrm{Heis}}_{\mathrm{ch}}, \kappa_{\mathrm{BKM}},
\kappa_{\mathrm{fiber}}\}$ on the row $\mathsf{C}$ of the
$5\times 5$ stratification matrix evaluate at conformal weight
$\lambda$:
\[
\kappa_{\mathrm{cat}}^{\beta\gamma_\lambda}=
6\lambda^2-6\lambda+1,
\qquad
\kappa^{\mathrm{Hodge}}_{\mathrm{ch}}(\beta\gamma_\lambda)=
\rho_\lambda\cdot\kappa_{\mathrm{cat}},
\qquad
\kappa^{\mathrm{Heis}}_{\mathrm{ch}}(\beta\gamma_\lambda)=0
\]
(no Heisenberg sub-current at generic $\lambda$),
$\kappa_{\mathrm{BKM}}=0$ and $\kappa_{\mathrm{fiber}}=0$.
\emph{Type signature:} (Open quadrant; chart-level scalar tuple;
Beilinson level~$5$; HP = five distinct measurements named;
collapse pattern $r_{\max}=4$ with quartic shadow $S_4=-5/12$).
The contact/quartic class-$\mathsf{C}$ collapse pattern is the
classification axis -- not the value of any single $\kappa$
(five-invariant discipline).
\end{remark}

\begin{remark}[Endpoint hypotheses for the $\beta\gamma$ surface]
\label{rem:betagamma-endpoint-hypotheses}
\index{endpoint hypothesis!beta-gamma@$\beta\gamma$}
Several quantitative results recorded for $\beta\gamma_\lambda$
in this chapter promote a finite-spin / quadratic / classical input
to a chiral statement; each carries an endpoint hypothesis package
(quadratic-candidate discipline):
\begin{enumerate}[label=\textup{(\arabic*)}]
\item \emph{Quartic contact coefficient} $S_4=-5/12$:
holds on the standard conformal-weight family
(Theorem~\ref{thm:betagamma-quartic-birth}), conditional on the
PBW-Koszul presentation existing for the bar complex through
degree~$4$; the rank-one weight-changing slice has $\mu=0$
by abelian rigidity (Theorem~\ref{thm:betagamma-rank-one-rigidity}).
\item \emph{Bar cohomology growth rate} $\sqrt{(1+x)/(1-3x)}$:
the discriminant $\Delta(x)=(1-3x)(1+x)$ shared with
$\widehat{\mathfrak{sl}}_2$ holds conditional on the lowest-weight
generator-type quotient definition, not the full bar group
(Theorem~\ref{thm:betagamma-bar-dim}).
\item \emph{Symplectic boson reduction} ($\lambda=1/2$):
$\kappa(\beta\gamma_{1/2})=-1/2$ and the $\mathbb{Z}_2$-equivariant
bar cohomology
(Proposition~\ref{prop:symplectic-equivariant-cohomology})
conditional on the symplectic-form preservation under the
weight-symmetric specialisation; this is not a Koszul self-duality
of $\beta\gamma_{1/2}$.
\item \emph{Vortex line $r$-matrix}: the explicit line $r$-matrix
(\S\ref{sec:betagamma-vortex-lines}) emerges from the regular contact
operator $\Theta^{\mathrm{ord}}_{\beta\gamma}=
\beta\otimes\gamma-\gamma\otimes\beta$, conditional on the
ordered bar transfer retaining the regular contact tensor
(this is not a pole-valued collision residue).
\end{enumerate}
None of (1)--(4) is asserted as an unconditional finite-spin /
quadratic / classical theorem.
\end{remark}

\section{The
\texorpdfstring{$\beta\gamma$}{beta-gamma} subtlety: three
invariants disagree maximally}
\label{sec:beta-gamma-subtlety}
\index{beta-gamma@$\beta\gamma$!three invariants disagree|textbf}
\index{three invariants!beta-gamma@$\beta\gamma$ as maximal case}

The $\beta\gamma$ system illustrates, in its sharpest form, the
importance of distinguishing the three invariants $p_{\max}$,
$k_{\max}$, $r_{\max}$ of
Chapter~\ref{ch:three-invariants}. The generator OPE
\[
\beta(z)\gamma(w) \;\sim\; \frac{1}{z-w},
\qquad
\beta(z)\beta(w) \;\sim\; 0,
\qquad
\gamma(z)\gamma(w) \;\sim\; 0
\]
contains a \emph{simple} pole only. Consequently:
\begin{itemize}
\item $p_{\max}(\beta\gamma) = 1$ (the generator OPE pole order);
\item $k_{\max}(\beta\gamma) = 0$ (the collision depth,
 $d\log$ absorption);
\item $r_{\max}(\beta\gamma) = 4$ (the shadow depth).
\end{itemize}
The discrepancy between $k_{\max} = 0$ and $r_{\max} = 4$ arises
because the shadow obstruction tower at higher degrees is \emph{not}
determined by the generator OPE alone. The composite current
$:\!\beta\gamma\!:$, a weight-$1$ field, carries a nontrivial quartic
contact class that places $\beta\gamma$ into shadow class~C, not into
the trivial class at $k_{\max} = 0$. The $(\beta\gamma)$-system is
thus the maximal case: the gap $r_{\max} - k_{\max} = 4$ is the
largest in the standard landscape at every fixed $p_{\max}$.

\begin{remark}[Pole order, collision depth, and shadow depth]
\label{warn:bg-depth-gap}
For $\beta\gamma_\lambda$ the OPE pole order, the pole-valued
collision residue, and the shadow tower are three different
invariants. The generator OPE has $p_{\max}=1$; the $d\log$
bar kernel lowers the pole-valued collision residue to
$k_{\max}=0$; the standard conformal-weight family has
$r_{\max}=4$ because the transferred contact coefficient
$S_4=-5/12$ survives.
\end{remark}

\begin{remark}[Implication for GZ\textup{26} commuting Hamiltonians]
\label{rem:bg-gz26-trivial}
\index{GZ26 commuting Hamiltonians!beta-gamma@$\beta\gamma$ trivial}
\index{beta-gamma@$\beta\gamma$!GZ26 Hamiltonians trivial}
By Theorem~\ref{thm:shadow-depth-operator-order}, the
GZ\textup{26} commuting Hamiltonians for the $\beta\gamma$ system are
trivial: they lie in the $k_{\max}=0$ trichotomy class. Thus
$\beta\gamma_\lambda$ has nontrivial finite shadow depth
$(r_{\max}=4)$ and trivial pole-valued commuting Hamiltonians at the
same time. The distinction is invisible from the central charge alone:
it depends on the triple
\[
p_{\max}=1,\qquad k_{\max}=0,\qquad r_{\max}=4 .
\]
\end{remark}

\section{Setup and conventions}
\index{beta-gamma system@$\beta\gamma$ system|textbf}

The $\beta\gamma$ system is the simplest bosonic chiral algebra
with two generators and a single simple-pole OPE channel.

\subsection{Algebraic structure}

Fields: $\beta(z)$ of conformal weight $h_\beta = \lambda$,
$\gamma(z)$ of weight $h_\gamma = 1 - \lambda$. The parameter
$\lambda$ always names the $\beta$-weight; the $\gamma$-weight is the
Serre-dual weight $1-\lambda$.

OPE:
\[
\beta(z)\gamma(w) = \frac{1}{z-w} + \text{regular},
\qquad
\beta(z)\beta(w) = \text{regular}, \quad \gamma(z)\gamma(w) = \text{regular}.
\]

\subsection{Mode algebra and stress tensor}

\begin{proposition}[Mode algebra \cite{FBZ04}; \ClaimStatusProvedElsewhere]
\label{prop:beta-gamma-modes}
Expand in modes:
\begin{align}
\beta(z) &= \sum_{n \in \mathbb{Z}} \beta_n z^{-n-\lambda},\qquad
\gamma(z) = \sum_{n \in \mathbb{Z}} \gamma_n z^{-n-(1-\lambda)}.
\end{align}
The commutation relations are $[\beta_m, \gamma_n] = \delta_{m+n,0}$.
\end{proposition}

\begin{theorem}[Stress tensor and central charge \cite{FBZ04}; \ClaimStatusProvedElsewhere]
\label{thm:beta-gamma-stress}
The stress-energy tensor is:
\[
T^{\beta\gamma}(z) = (1-\lambda) \normord{\beta(z)\partial\gamma(z)} - \lambda \normord{\partial\beta(z)\gamma(z)}
\]
This generates the Virasoro algebra with central charge:
\[
c_{\beta\gamma} = +2(6\lambda^2 - 6\lambda + 1) \quad \text{(bosonic $\beta\gamma$)}
\]
For the \emph{fermionic} $bc$ system with the same conformal weights, the sign reverses:
$c_{bc} = 1 - 3(2\lambda - 1)^2 = -2(6\lambda^2 - 6\lambda + 1)$.
\end{theorem}

\begin{computation}[Central charges for special cases]
\label{comp:beta-gamma-central-charges}
\label{thm:beta-gamma-ope-complete}
Fermionic $bc$ system ($c_{bc} = 1 - 3(2\lambda-1)^2$):
$\lambda = 2$: $c = -26$ (string theory ghosts);
$\lambda = 1$: $c = -2$;
$\lambda = 1/2$: $c = 1$ (complex fermion pair);
$\lambda = 0$: $c = -2$.
Bosonic $\beta\gamma$ system ($c_{\beta\gamma} = +2(6\lambda^2 - 6\lambda + 1)$):
$\lambda = 1$: $c = +2$;
$\lambda = 0$: $c = +2$;
$\lambda = 1/2$: $c = -1$ (symplectic bosons).
\end{computation}

\section{Bar complex computation}

The bar complex of $\beta\gamma$ computes the degree-by-degree
projections of the universal MC element
$\Theta_{\beta\gamma} := D_{\beta\gamma} - d_0 \in \MC(\gAmod)$
(Theorem~\ref{thm:mc2-bar-intrinsic}). At degree~$2$,
the degree-$2$ shadow restricted to the weight-changing line vanishes
($\Theta_{\beta\gamma}^{\leq 2}\big|_{\mathrm{w.c.}} = 0$:
the $\beta\gamma$ OPE has no double pole, so
the bar complex is uncurved on this slice).
The global modular characteristic
$\kappa(\beta\gamma) = c/2 \neq 0$
(Proposition~\ref{prop:betagamma-obstruction-coefficient})
is nonzero; the vanishing is specific to the weight-changing
direction. The mixed OPE $\beta\gamma\sim 1/(z{-}w)$ contributes the
bar differential, but the class-$\mathsf C$ scalar cubic coefficient
on the standard conformal-weight family is $S_3=0$. The quartic
contact coefficient is $S_4=-5/12$ on that standard family, while
its rank-one weight-line scalar $\mu_{\beta\gamma}$ is zero. The
quintic and higher coefficients vanish on the class-$\mathsf C$
family, so the standard-family shadow tower terminates at degree~$4$.

\medskip
Unlike the Heisenberg bar complex (\S\ref{sec:frame-bar-all}), where
$d_{\mathrm{bracket}} = 0$ because the abelian OPE has no simple pole,
$\beta\gamma$ has a nontrivial bar differential: the mixed OPE
$\beta(z)\gamma(w) \sim 1/(z-w)$ contracts mixed ordered-contact
pairs by extracting the coefficient of the simple pole. The symbol
$\eta_{ij}=d\log(z_i-z_j)$ records the boundary orientation in the
ordered bar complex; it is not multiplied with the OPE to form a
literal double-pole residue. Thus
\[
 \operatorname{Coeff}_{(z_i-z_j)^{-1}}
 \bigl(\beta(z_i)\gamma(z_j)\bigr)=1
\]
is the chain-level contraction,
while pure $\beta^{\otimes n}$ and $\gamma^{\otimes n}$ tensors are
cycles until higher bar boundaries are imposed. This ordered contact
extraction is the chain-level operation; it is not the closed
collision-residue $r$-matrix, which is zero after $d\log$
absorption. This is the mechanism behind the growth of the
Koszul-dual coalgebra; the bare bar complex is not the acyclic object.

\subsection{Degree by degree analysis}

\begin{theorem}[Lowest-weight bar skeleton through degree~3;
\ClaimStatusConditional]
\label{thm:betagamma-complete-bar}
The lowest-weight generator-type skeleton of the $\beta\gamma$ bar
complex through degree~$3$ is as follows. The unreduced bar group is
infinite-dimensional in every positive bar degree because the
augmentation ideal contains all conformal descendants; the displayed
complex keeps the two primary types and the first derivative sectors
created by the simple-pole coefficient extraction.

\emph{Degree~0.} $\bar{B}^0 = \mathbb{C}|0\rangle$ (vacuum)

\emph{Degree~1.} $\bar{B}^1 = V_\beta \oplus V_\gamma$ where
\[
V_\beta = \operatorname{span}\{\beta_{-n-h_\beta}|0\rangle : n \geq 0\},
\qquad
V_\gamma = \operatorname{span}\{\gamma_{-n-h_\gamma}|0\rangle : n \geq 0\}.
\]

\emph{Degree~2.}
\begin{align}
\bar{B}^2 = &(V_\beta \otimes V_\beta) \oplus (V_\gamma \otimes V_\gamma) \\
&\oplus (V_\beta \otimes V_\gamma) \oplus (V_\gamma \otimes V_\beta) \\
&\oplus V_{\partial\beta} \oplus V_{\partial\gamma}
\end{align}

The differential $d: \bar{B}^2 \to \bar{B}^1$:
\begin{align}
d(\beta \otimes \beta) &= 0 \text{ (no pole in OPE)} \\
d(\gamma \otimes \gamma) &= 0 \\
d(\beta \otimes \gamma)
&= \operatorname{Coeff}_{(z_1-z_2)^{-1}}
\bigl(\beta(z_1)\gamma(z_2)\bigr) = 1 \\
d(\gamma \otimes \beta) &= -1 \\
d(\partial\beta) &= 0, \quad d(\partial\gamma) = 0
\end{align}

\emph{Degree~3.} Number of type components $= 2 \cdot 3^2 = 18$

Pure tensor components ($2^3 = 8$ terms): $(V_\beta)^{\otimes 3}$ and $(V_\gamma)^{\otimes 3}$ (each 1-dimensional), $(V_\beta)^{\otimes 2} \otimes V_\gamma$ and $V_\beta \otimes (V_\gamma)^{\otimes 2}$ (each with 3 orderings). The remaining 10 derivative types arise from replacing one or both of positions 2, 3 by $\partial\gamma$.

Key differential:
\[d(\beta_1 \otimes \beta_2 \otimes \gamma_3) = \beta_1 \otimes 1 - 1 \otimes \beta_2\]

The rank (number of independent generator types) satisfies $\operatorname{rank}(\bar{B}^n) = 2 \cdot 3^{n-1}$ for all $n \geq 1$ (Theorem~\ref{thm:betagamma-bar-dim}).
\end{theorem}

\begin{proof}
The explicit computation for degrees~1--3 proceeds as follows.

\emph{Degree~1:} $\bar{B}^1 = V_\beta \oplus V_\gamma$ (2 types).

\emph{Degree~2:} There are $2^2 = 4$ pure tensor types (choosing $\beta$ or $\gamma$ at each of 2 positions) plus 2 derivative types ($V_{\partial\beta}$ and $V_{\partial\gamma}$), giving $4 + 2 = 6$ types.

\emph{Degree~3:} There are $2^3 = 8$ pure tensor types. Derivative terms arise when one or both of positions~2, 3 are replaced by conformal descendants: substituting $\partial\gamma$ at position~2 (with any of $\beta, \gamma$ at positions~1, 3) gives $2 \cdot 2 = 4$ types, substituting at position~3 gives another~4, and substituting at both positions~2 and~3 gives $2$, for a total of $8 + 4 + 4 + 2 = 18$ types. These match $2 \cdot 3^0 = 2$, $2 \cdot 3^1 = 6$, $2 \cdot 3^2 = 18$.
\end{proof}

\subsection{Cohomology calculation}

\begin{theorem}[Bar cohomology of \texorpdfstring{$\beta\gamma$}{beta-gamma}; \ClaimStatusConditional]
\label{thm:betagamma-bar-cohomology}
For generic conformal weight $\lambda$, the geometric bar cohomology
of the $\beta\gamma$ system is the Koszul-dual coalgebra
\[
H^\bullet\!\bigl(\bar{B}_{\mathrm{geom}}(\beta\gamma_\lambda)\bigr)
\;\cong\;(\beta\gamma_\lambda)^{\mathrm{i}},
\qquad
\bigl((\beta\gamma_\lambda)^{\mathrm{i}}\bigr)^\vee
\;\cong\;
(\beta\gamma_\lambda)^!
\;\simeq\;
bc_\lambda .
\]
It is not concentrated in degree~$0$. The acyclic object is the
twisted Koszul resolution
$\beta\gamma_\lambda\otimes_{\tau}\bar{B}(\beta\gamma_\lambda)$,
whose augmentation has cohomology $\mathbb{C}$.
\end{theorem}

\begin{proof}
Theorem~\ref{thm:betagamma-bc-koszul} identifies the Verdier/Koszul
dual algebra with $bc_\lambda$. Equivalently, the orthogonal
complement of the antisymmetric $\beta\gamma$ relation is the
Clifford relation for $bc_\lambda$. The bar construction computes the
dual coalgebra before finite dualization:
$H^\bullet(\bar B(A))=A^{\mathrm i}$ and
$A^!=(A^{\mathrm i})^\vee$ on the finite-type truncations.
Thus the positive bar cohomology classes are precisely the coalgebra
generators and relations dual to the $bc$ system.

The degree-$2$ calculation below already exhibits the distinction:
$[\beta\otimes\beta]$ and $[\gamma\otimes\gamma]$ survive as
Koszul-dual coalgebra classes in $\bar B(\beta\gamma)$, whereas mixed
classes are contracted by the simple-pole coefficient of
$\beta(z)\gamma(w)\sim 1/(z-w)$. Exactness belongs to the twisted
tensor product
$\beta\gamma_\lambda\otimes_\tau\bar B(\beta\gamma_\lambda)$, the
Koszul resolution of the vacuum, not to the bare bar coalgebra.
\end{proof}

\section{Koszul dual}
\label{sec:betagamma-koszul-dual}

\subsection{Main result: Koszul duality with free fermions}

\begin{theorem}[Koszul dual of \texorpdfstring{$\beta\gamma$}{beta-gamma}; \ClaimStatusProvedHere]
\label{thm:betagamma-fermion-koszul}
\index{Koszul duality!beta-gamma vs.\ free fermion@$\beta\gamma$ vs.\ free fermion|textbf}
The Koszul dual of the $\beta\gamma$ system is the $bc$ ghost system:
\[(\beta\gamma)^! \cong bc \quad \text{and} \quad (bc)^! \cong \beta\gamma\]

This is a \emph{two-generator} duality ($\dim V = 2$: one bosonic pair
$\beta, \gamma$ dualizes to one fermionic pair $b, c$), not to be confused with
the single-generator duality $\Lambda^{\mathrm{ch}}(\psi) \leftrightarrow
\mathrm{Sym}^{\mathrm{ch}}(\gamma)$. The $bc$ system has:
\begin{enumerate}
\item \emph{Generators.} $b$ with conformal weight $h_b = \lambda$ and $c$ with conformal weight $h_c = 1 - \lambda$
\item \emph{Statistics.} Fermionic ($b$ and $c$ are odd/anticommuting)
\item \emph{OPE.} $b(z)c(w) \sim \dfrac{1}{z-w}$, \quad $b(z)b(w) \sim 0$, \quad $c(z)c(w) \sim 0$
\end{enumerate}
\end{theorem}

\begin{proof}
This is the chiral analog of the classical Koszul duality between $\text{Sym}(V)$ and $\Lambda(V^*)$
for associative algebras. The proof follows from the Gui--Li--Zeng framework~\cite{GLZ22}.

\emph{Step~1: Quadratic Data}

For $\beta\gamma$ system:
\begin{itemize}
\item Generators: $N = \mathcal{O}_X \cdot \beta \oplus \mathcal{O}_X \cdot \gamma$
\item Relations: $P \subset j_* j^*(N \boxtimes N)$ encode the ordered Weyl pairing:
\[
\Theta_{\beta\gamma}^{\mathrm{ord}}
=\beta \boxtimes \gamma-\gamma \boxtimes \beta,
\qquad
\langle \beta, \gamma \rangle = \frac{1}{z_1 - z_2}, \quad
\langle \gamma, \beta \rangle = -\frac{1}{z_1-z_2}.
\]
The pure pairings $\langle \beta,\beta\rangle$ and
$\langle\gamma,\gamma\rangle$ vanish.
\end{itemize}

For the $bc$ ghost system $\mathcal{F}$:
\begin{itemize}
\item Generators: $N' = \mathcal{O}_X \cdot b \oplus \mathcal{O}_X \cdot c$ (two fermionic generators, dual to $\beta, \gamma$)
\item Relations: $P' \subset j_* j^*(N' \boxtimes N')$ encode the graded Clifford pairing:
\[
\Theta_{bc}^{\mathrm{ord}}
=b \boxtimes c+c \boxtimes b,
\qquad
b \boxtimes b=0,\quad c\boxtimes c=0.
\]
The plus sign is the graded tensor sign for odd generators; the
ordinary desuspended tensor representative carries the corresponding
Koszul minus sign.
\end{itemize}

\emph{Step~2: Dual Quadratic Data}

The Koszul dual is constructed via:
\[(s^{-1}N^{\vee}\omega^{-1}, P^{\perp})\]

For $\mathcal{F}$: The dual of the graded Clifford relation
$b \boxtimes c + c \boxtimes b$ is the Weyl/symplectic pairing
$\beta \boxtimes \gamma - \gamma \boxtimes \beta$, recovering
$\beta\gamma$.

For $\beta\gamma$: The dual of the Weyl pairing is the Clifford relation, recovering $\mathcal{F}$.
\end{proof}

\subsection{Bar-cobar verification}
\label{subsec:bar-cobar-verification}

\subsubsection{\texorpdfstring{Bar complex of the $bc$ ghost system}{Bar complex of the bc ghost system}}

\begin{proposition}[Bar complex structure; \ClaimStatusProvedHere]
\label{prop:bar-bc-system}
The bar complex of the $bc$ ghost system $\mathcal{F}$ has the structure:
\[\bar{B}^n(\mathcal{F}) = \mathrm{Sym}^n_c(s^{-1}\bar{\mathcal{F}})\]
where $\bar{\mathcal{F}} = \ker(\epsilon)$ is the augmentation ideal spanned by modes of both $b$ and $c$.

\emph{Explicit computation} (at lowest conformal weights):
\begin{align}
\bar{B}^0(\mathcal{F}) &= \mathbb{C}|0\rangle \\
\bar{B}^1(\mathcal{F}) &= \mathbb{C}\langle b_{-\lambda}|0\rangle, c_{-(1-\lambda)}|0\rangle, \ldots \rangle \\
\bar{B}^2(\mathcal{F}) &= \mathbb{C}\langle b \otimes b,\; c \otimes c,\; b \otimes c,\; c \otimes b \rangle_{\text{modes}}
\end{align}

The symmetric coalgebra structure arises because $\mathcal{F}$ is a
Clifford-type exterior algebra. By Koszul duality, the bar complex
of an exterior algebra resolves the symmetric Koszul-dual coalgebra:
\[
H^*(\bar{B}(\Lambda(V))) \cong \mathrm{Sym}^c(V^*),
\]
and the Koszul-dual algebra is
$\Lambda(V)^! \cong \mathrm{Sym}(V^*)$ under finite-type duality.
This is the chiral analogue for the $bc$ system.
\end{proposition}

\begin{proof}
The $bc$ ghost system has OPEs $b(z)c(w) \sim \frac{1}{z-w}$,
$b(z)b(w) \sim 0$, $c(z)c(w) \sim 0$, with graded Clifford
ordered-contact tensor
\[
R_{\mathrm{Cliff}}
=b(z_1)\otimes c(z_2)+c(z_2)\otimes b(z_1).
\]
The bar complex $\bar{B}(\mathcal{F})$ is the cofree coalgebra on the desuspended augmentation ideal with differential encoding the OPE.

The ordered differential $d_{\mathrm{ord}}$ on $\bar{B}^2$ extracts
the simple-pole coefficient of the OPE at collision:
\[
d_{\mathrm{ord}}(b \otimes c)
=
\operatorname{Coeff}_{(z_1-z_2)^{-1}}
\bigl(b(z_1)c(z_2)\bigr)=1,
\]
mapping to $\bar{B}^0$. Similarly
$d_{\mathrm{ord}}(c \otimes b) = -1$. This minus is the ordinary
ordered-tensor sign after desuspension; it does not replace
the graded Clifford tensor $b\otimes c+c\otimes b$. Meanwhile
$d_{\mathrm{ord}}(b \otimes b) = 0$ and
$d_{\mathrm{ord}}(c \otimes c) = 0$ since $b(z)b(w)$ and
$c(z)c(w)$ are regular.
Hence $H^2(\bar{B}(\mathcal{F}))
= \mathbb{C}\langle [b \otimes b], [c \otimes c] \rangle$,
whose dual gives the symmetric (commutative) relations of the $\beta\gamma$ system.
\end{proof}

\subsubsection{\texorpdfstring{Koszul dual algebra: $\mathcal{F}^! = H^\bullet\bar{B}(\mathcal{F})^\vee \cong \beta\gamma$}{Koszul dual algebra: F! = H(B(F))v = beta-gamma}}

\begin{theorem}[Koszul dual of the free fermion; \ClaimStatusProvedHere]
\label{thm:cobar-betagamma}
The Koszul dual algebra of the free fermion is the $\beta\gamma$ system:
\[\mathcal{F}^! =
\bigl(H^\bullet\bar{B}(\mathcal{F})\bigr)^\vee
\cong \text{Chiral algebra}(\beta, \gamma \mid [\beta,\gamma] = 1).\]
(Bar-cobar inversion gives $\Omega(\bar{B}(\mathcal{F})) \cong \mathcal{F}$, recovering the fermion itself.)
\end{theorem}

\begin{proof}
\emph{Step~1.} Dualize $\bar{B}^1(\mathcal{F})$ to get two generators $\beta, \gamma$.

\emph{Step~2.} The graded Clifford relation in
$\bar{B}^2(\mathcal{F})$ dualizes to the ordinary Weyl tensor:
\[\beta \boxtimes \gamma - \gamma \boxtimes \beta = \frac{1}{z_1-z_2} \cdot 1\]

This is the symplectic pairing.

\emph{Step~3.} The dual coalgebra structure encodes the chiral product:
\[\beta(z)\gamma(w) = \frac{1}{z-w} + \text{regular}\]
\end{proof}

\subsubsection{Bar complex of beta-gamma (detailed)}

\begin{proposition}[\ClaimStatusProvedHere]
\label{prop:betagamma-bar-deg2}
The bar complex $\bar{B}(\beta\gamma)$ has the following structure:

\emph{Degree~2 cohomology.}
\[H^2(\bar{B}(\beta\gamma)) = \mathbb{C}\langle [\beta \otimes \beta], [\gamma \otimes \gamma] \rangle\]

where both classes are represented by cycles that are \emph{not} boundaries (since $d^3$ maps into the span of elements involving $\beta \otimes \gamma$ or $\gamma \otimes \beta$, not pure types). The acyclicity of mixed terms follows from the relation:
\[d(\beta \otimes \gamma) = 1\]
\end{proposition}

\begin{proof}
We compute:
\begin{align}
d(\beta \otimes \beta) &= 0 \quad \text{(no pole)} \\
d(\gamma \otimes \gamma) &= 0 \quad \text{(no pole)} \\
d(\beta \otimes \gamma)
&=
\operatorname{Coeff}_{(z_1-z_2)^{-1}}
\bigl(\beta(z_1)\gamma(z_2)\bigr)=1 \\
d(\gamma \otimes \beta) &= -1
\end{align}

Therefore in cohomology:
\[[\beta \otimes \beta] \neq 0, \quad [\gamma \otimes \gamma] \neq 0\]
but these satisfy:
\[[\beta \otimes \beta] \cdot [\gamma] = [\gamma \otimes \gamma] \cdot [\beta] = 0\]
\end{proof}

\subsubsection{\texorpdfstring{Koszul dual: $(\beta\gamma)^! = H^\bullet\bar{B}(\beta\gamma)^\vee \cong \mathcal{F}$}{Koszul dual: (bg)! = F}}

\begin{theorem}[Koszul dual of \texorpdfstring{$\beta\gamma$}{beta-gamma}; \ClaimStatusProvedHere]
\label{thm:cobar-fermions}
The Koszul dual algebra of the $\beta\gamma$ system is the free fermion ($bc$ ghost system):
\[(\beta\gamma)^! =
\bigl(H^\bullet\bar{B}(\beta\gamma)\bigr)^\vee
\cong \text{Chiral algebra}(b, c \mid b^2 = 0,\; c^2 = 0,\; b(z)c(w) \sim \tfrac{1}{z-w}).\]
(Bar-cobar inversion gives $\Omega(\bar{B}(\beta\gamma)) \cong \beta\gamma$, recovering $\beta\gamma$ itself.)
\end{theorem}

\begin{proof}
\emph{Step~1.} Dualize $\bar{B}^1(\beta\gamma) = \mathbb{C}\langle \beta \rangle \oplus \mathbb{C}\langle \gamma \rangle$
to get two fermionic generators $b = \beta^*$, $c = \gamma^*$.

\emph{Step~2.} The exterior coalgebra structure (from $\beta \cdot \beta = \gamma \cdot \gamma = 0$) dualizes to the anticommutation relations:
\[b^2 = 0, \quad c^2 = 0, \quad b(z)c(w) \sim \frac{1}{z-w}\]

\emph{Step~3.} This defines the $bc$ ghost system (free fermions with two generators).
\end{proof}

\subsection{Geometric interpretation}
\label{subsec:NAP-betagamma}

\begin{remark}[Verdier duality]
Verdier duality on configuration spaces exchanges exterior algebras (antisymmetric pairing, fermionic) with polynomial algebras (symplectic pairing, bosonic); see Beilinson--Drinfeld~\cite{BD04} \S3.8.
\end{remark}

\section{Relationship to special cases}

\subsection{\texorpdfstring{Understanding the $\lambda$ parameter}{Understanding the parameter}}

\begin{remark}[Conformal weights]
The $\beta\gamma$ system has fields with conformal weights:
\[h_\beta = \lambda, \quad h_\gamma = 1-\lambda\]
where $\lambda$ is a parameter.

When $\lambda = 1$, $\beta$ has weight~1 and $\gamma$ has weight~0. The system remains \emph{bosonic} (commuting fields) for all values of~$\lambda$. The $bc$ ghost system is the fermionic counterpart with anticommuting fields of the same weights.
\end{remark}

\subsection{\texorpdfstring{The $bc$ ghost system and Koszul duality}{The bc ghost system and Koszul duality}}

\begin{proposition}[Central charge complementarity for \texorpdfstring{$\beta\gamma$}{beta-gamma}/\texorpdfstring{$bc$}{bc}; \ClaimStatusProvedHere]\label{prop:betagamma-bc-koszul-detailed}
\index{bc ghost system@$bc$ ghost system|textbf}
\index{statistics exchange!bosonic vs.\ fermionic}
\index{central charge!complementarity}
The Koszul duality $(\beta\gamma)^! \simeq bc$ of Theorem~\ref{thm:betagamma-fermion-koszul}
has a sharp central charge consequence:
\begin{equation}\label{eq:bg-bc-koszul}
c_{\beta\gamma} = +2(6\lambda^2 - 6\lambda + 1),
\qquad
c_{bc} = 1 - 3(2\lambda - 1)^2 = -2(6\lambda^2 - 6\lambda + 1),
\qquad
c_{\beta\gamma} + c_{bc} = 0.
\end{equation}
\end{proposition}

\begin{proof}
The central charge of the $\beta\gamma$ system is computed from the stress tensor (Theorem~\ref{thm:beta-gamma-stress}). Replacing bosonic statistics by fermionic statistics negates the Virasoro central term: the commutator $[L_m, L_n]$ picks up a sign from the fermion loop, so $c_{bc} = -c_{\beta\gamma}$.

This exact cancellation is specific to the Sym/Ext exchange.
For general Koszul dual pairs the central charges need not sum to zero:
affine Kac--Moody algebras satisfy
$c(k) + c(-k-2h^\vee) = c_{\mathrm{crit}}$ instead
(Remark~\ref{rem:km-central-charge-sum}).
\end{proof}

\begin{remark}[Statistics exchange as a general principle]\label{rem:statistics-exchange}
The $\beta\gamma \leftrightarrow bc$ duality illustrates a general principle: chiral Koszul duality exchanges bosonic and fermionic statistics. At the operadic level, this reflects the classical duality $\mathrm{Com}^! = \mathrm{Lie}$ (commutative $\leftrightarrow$ Lie), which at the level of enveloping algebras becomes $\mathrm{Sym}(V)^! = \Lambda(V^*)$ (symmetric $\leftrightarrow$ exterior). The chiral lift preserves this exchange because the quadratic duality functor $(N, P) \mapsto (s^{-1}N^\vee\omega^{-1}, P^\perp)$ sends symmetric pairings to exterior relations and vice versa.
\end{remark}

\subsection{Boson-fermion correspondence}

\begin{theorem}[Physical bosonization \cite{FBZ04}; \ClaimStatusProvedElsewhere]
\label{thm:physical-bosonization}
There is a physics ``bosonization'' map relating $\beta\gamma$ and free fermions at the level of
\emph{correlation functions}:
\[\text{Correlators}[\beta\gamma] \xleftrightarrow{\text{Bosonization}} \text{Correlators}[\mathcal{F}]\]

This is \emph{different} from Koszul duality: bosonization equates correlation functions, while Koszul duality is an algebraic relationship between chiral algebras via bar-cobar
(Remark~\ref{rem:bosonization-not-koszul}).
\end{theorem}

\subsection{\texorpdfstring{Symplectic bosons ($\lambda = 1/2$)}{Symplectic bosons ( = 1/2)}}

At $\lambda = 1/2$, both fields have weight $1/2$:
\[T = \frac{1}{2}(\beta\partial\gamma - \partial\beta \cdot \gamma)\]

At this value the local OPE with the stress tensor is still the
ordinary primary-field OPE. Logarithmic phenomena occur in the module
category and in equivariant/orbifold constructions, not as a
logarithmic singularity in $T(z)\beta(w)$ or $T(z)\gamma(w)$. The
Koszul dual remains the fermionic $bc_{1/2}$ system.

\section{Geometric realization}

\subsection{Bundle data on curves}

\begin{construction}[Geometric \texorpdfstring{$\beta$}{beta}-\texorpdfstring{$\gamma$}{gamma} system]
\label{const:geometric-beta-gamma}
\index{canonical bundle!in beta-gamma realization@in $\beta\gamma$ realization}
On a curve $X$, the $\beta$-$\gamma$ system with parameter $\lambda$ is geometrically:
\begin{align}
\beta &\in \Gamma(X, K_X^\lambda \otimes \mathcal{L})\\
\gamma &\in \Gamma(X, K_X^{1-\lambda} \otimes \mathcal{L}^*)
\end{align}

where $K_X$ is the canonical bundle and $\mathcal{L}$ is an auxiliary line bundle.

\emph{Special cases.}
\begin{enumerate}
\item $\lambda = 1$: $\beta \in \Gamma(K_X)$ (differentials), $\gamma \in \Gamma(\mathcal{O}_X)$ (functions)
\item $\lambda = 2$: $\beta \in \Gamma(K_X^2)$ (quadratic differentials), $\gamma \in \Gamma(K_X^{-1})$ (vector fields); the fermionic counterpart at $\lambda = 2$ provides the BRST ghosts of the bosonic string
\item $\lambda = 1/2$: Both fields are sections of $K_X^{1/2}$ (spin structures required)
\end{enumerate}
\end{construction}

\begin{remark}[Bundle pairing as bar differential]
\label{rem:bundle-bar-mechanism}
The Serre duality residue pairing
$K_X^\lambda \otimes K_X^{1-\lambda} \to K_X
\xrightarrow{\mathrm{Res}} \mathbb{C}$ is precisely the geometric
source of the ordered bar differential:
\[
 \operatorname{Coeff}_{(z-w)^{-1}}\beta(z)\gamma(w)=1,
 \qquad
 \operatorname{Coeff}_{(z-w)^{-1}}\gamma(z)\beta(w)=-1
\]
in the Weyl ordering. This extracts the contact coefficient along
the diagonal $\Delta \hookrightarrow X \times X$. Pure
$\beta$--$\beta$ or $\gamma$--$\gamma$ collisions produce no pole
because $K_X^\lambda \otimes K_X^\lambda$ has no canonical residue;
this is the geometric reason that only mixed types are killed by $d$,
giving the $3^n$ growth. This ordered contact extraction is separate
from the closed pole-valued collision residue
$r^{\mathrm{coll}}_{\beta\gamma}(z)=0$.
\end{remark}

\subsection{Configuration space picture}

The bar complex elements at type $(n,m)$ are sections over the FM compactification:
\[\omega_{n,m} \in \Gamma\bigl(\overline{C}_{n+m}(X),\; (K_X^\lambda)^{\boxtimes n} \boxtimes (K_X^{1-\lambda})^{\boxtimes m} \otimes \Omega^*_{\log}\bigr)\]
In coordinates:
\[\omega_{n,m} = \beta(z_1) \cdots \beta(z_n) \gamma(w_1) \cdots \gamma(w_m) \prod_{i<j} \eta_{ij}\]
The differential extracts the simple-pole coefficient only along
mixed collision divisors $D_{z_i = w_j}$:
\[
d(\omega_{n,m})
= \sum_{i,j}
\operatorname{Coeff}_{(z_i-w_j)^{-1}}[\omega_{n,m}]
= \sum_{i,j}\omega_{n-1,m-1}\big|_{z_i=w_j}.
\]
Thus the algebraic bar differential is geometrically the coefficient
map along the simple-pole locus of the Serre pairing; the logarithmic
form supplies the orientation of the boundary stratum.

\subsection{Geometric bar complex}

\begin{theorem}[Bar complex of the \texorpdfstring{$\beta$}{beta}-\texorpdfstring{$\gamma$}{gamma} system; \ClaimStatusProvedHere]
\label{thm:beta-gamma-bar}
The bar complex of the $\beta$-$\gamma$ system is:
\[
\bar{B}^n = \bigl(\mathrm{Free}[\beta, \gamma]\bigr)^{\otimes (n+1)} \otimes \Omega^n(\overline{C}_{n+1}(X))
\]
The differential is ordered-contact coefficient extraction:
\[
d = \sum_{i<j}
\operatorname{Coeff}_{(z_i-z_j)^{-1}}
\bigl[\Theta_{\beta\gamma,ij}^{\mathrm{ord}}\bigr]\,
\operatorname{Res}^{\log}_{D_{ij}}(\eta_{ij}\wedge -),
\qquad
\Theta_{\beta\gamma,ij}^{\mathrm{ord}}
=\beta_i\otimes\gamma_j-\gamma_i\otimes\beta_j,
\]
where $\eta_{ij} = dz_i/(z_i-z_j)$ are logarithmic forms.
Here $\operatorname{Res}^{\log}_{D_{ij}}(\eta_{ij})=1$ is the
logarithmic boundary residue of the form alone; the singular OPE
contributes only through the separate coefficient extraction.
\end{theorem}

\begin{proof}
This is the specialization of the geometric bar construction
(Definition~\ref{def:geometric-bar}) to the $\beta\gamma$ system.
The bar degree $n$ component uses $n+1$ copies of the free field algebra
on the FM compactification $\overline{C}_{n+1}(X)$, with the differential
given by iterated coefficient extraction along collision divisors
$D_{ij} = \{z_i = z_j\}$.
The $\beta\gamma$ OPE $\beta(z)\gamma(w) \sim 1/(z-w)$ determines the
coefficient
$\operatorname{Coeff}_{(z_i-z_j)^{-1}}
[\Theta_{\beta\gamma,ij}^{\mathrm{ord}}]$. The opposite ordered
channel carries the minus sign in
$\Theta_{\beta\gamma}^{\mathrm{ord}}$; the closed collision projection
still has $r^{\mathrm{coll}}_{\beta\gamma}(z)=0$ after the $d\log$
absorption convention.
\end{proof}

\subsection{Wakimoto connection and universal property}

\begin{remark}[Connection to Wakimoto realization]
The Wakimoto free field realization uses $\beta$-$\gamma$ systems extensively:
\[
\mathcal{M}_{\mathrm{Wak}} = \mathrm{Free}[\phi_i] \otimes \bigotimes_{\alpha \in \Delta_+} \mathrm{Free}[\beta_\alpha, \gamma_\alpha]
\]
Each root $\alpha$ contributes a $\beta$-$\gamma$ system. These are the building blocks
for the free field realization of affine Kac--Moody and $\mathcal{W}$-algebras.
\end{remark}

\begin{theorem}[Universal property of \texorpdfstring{$\beta$}{beta}-\texorpdfstring{$\gamma$}{gamma} \cite{FBZ04}; \ClaimStatusProvedElsewhere]
\label{thm:beta-gamma-universal}
The $\beta$-$\gamma$ system is the \emph{free vertex algebra} generated by two fields
$\beta, \gamma$ of weights $\lambda, 1-\lambda$ with the single relation
$\beta(z)\gamma(w) \sim 1/(z-w)$.
For any vertex algebra $V$ with fields $\beta', \gamma'$ satisfying
this OPE, there exists a unique homomorphism
$\mathrm{Free}[\beta, \gamma] \to V$
sending $\beta \mapsto \beta'$, $\gamma \mapsto \gamma'$.
\end{theorem}

%===================================================================================
% DEGREE-BY-DEGREE BAR DIFFERENTIAL: DEGREES 4 AND 5
%===================================================================================

\section{Degree-by-degree bar differential: degrees 4 and 5}
\label{sec:betagamma-bar-degree45}
\index{beta-gamma system@$\beta\gamma$ system!bar complex!degree 4 and 5}

The bar complex $\bar{B}(\beta\gamma)$ in degrees~4 and~5 extends the computation of \S\ref{subsec:bar-cobar-verification}.

\subsection{Type enumeration}

At bar degree~$n$, an element of $\bar{B}^n(\beta\gamma)$ is a section
over $\overline{C}_{n+1}(X)$ of the form
\[
a_1(z_1) \otimes \cdots \otimes a_n(z_n) \otimes \omega,
\qquad \omega \in \Omega^n_{\log}(\overline{C}_{n+1}(X)),
\]
where each $a_i$ is a monomial in $\beta, \gamma$ and their
conformal descendants $\partial^k \beta, \partial^k \gamma$.
A \emph{type} is specified by the choice of generator
($\beta$, $\gamma$, or a derivative $\partial^j \beta$,
$\partial^j \gamma$ with $j \geq 1$) at each tensor position.
At each position, there are three independent families of
generators: $\beta$, $\gamma$, and a conformal descendant
(the first derivative $\partial \beta$ or $\partial \gamma$
contributes at the leading order; higher derivatives contribute at
higher conformal weight but do not introduce new types at a given
weight truncation). Thus the type count satisfies:

\begin{equation}\label{eq:betagamma-type-count}
\text{types at degree } n = 2 \cdot 3^{n-1}, \qquad n \geq 1.
\end{equation}

\begin{computation}[Dimension table through degree~10]
\label{comp:betagamma-dim-table}
\index{beta-gamma system@$\beta\gamma$ system!bar complex!dimension table}
The number of generator types in $\bar{B}^n(\beta\gamma)$ through degree~10:
\begin{center}
\begin{tabular}{ccl}
\toprule
Degree $n$ & Types $= 2 \cdot 3^{n-1}$ & Description \\
\midrule
0 & 1 & vacuum $\mathbb{C}|0\rangle$ \\
1 & 2 & $\beta$, $\gamma$ \\
2 & 6 & $4$ pure $+ 2$ derivative \\
3 & 18 & $8$ pure $+ 10$ derivative \\
4 & 54 & $16$ pure $+ 38$ derivative \\
5 & 162 & $32$ pure $+ 130$ derivative \\
6 & 486 & $64$ pure $+ 422$ derivative \\
7 & 1458 & $128$ pure $+ 1330$ derivative \\
8 & 4374 & $256$ pure $+ 4118$ derivative \\
9 & 13122 & $512$ pure $+ 12610$ derivative \\
10 & 39366 & $1024$ pure $+ 38342$ derivative \\
\bottomrule
\end{tabular}
\end{center}
At degree~$n$, the $2^n$ pure types correspond to choosing $\beta$ or
$\gamma$ at each of $n$ positions; the remaining $2 \cdot 3^{n-1} - 2^n$
types involve at least one conformal descendant. The exponential growth
rate $3^n$ reflects the three generator types ($\beta$, $\gamma$, derivative)
at each tensor position.
\end{computation}

\subsection{Degree~4: explicit differentials}

At degree~4, the bar complex element is a section over $\overline{C}_5(X)$:
\[
a_1 \otimes a_2 \otimes a_3 \otimes a_4 \otimes \omega,
\qquad \omega \in \Omega^4_{\log}(\overline{C}_5(X)).
\]
The logarithmic $4$-forms are spanned by products
$\eta_{i_1 j_1} \wedge \eta_{i_2 j_2} \wedge \eta_{i_3 j_3} \wedge \eta_{i_4 j_4}$
modulo Arnold relations. By the Orlik--Solomon algebra
(Theorem~\ref{thm:arnold-relations-appendix}), the space
$\Omega^4_{\log}(\overline{C}_5)$ has dimension
$\binom{4}{4} \cdot 4! / \ldots$; the relevant basis is the set of
no-broken-circuit (NBC) monomials of degree~4 in the $\eta_{ij}$.

\begin{computation}[Degree-4 bar differential]
\label{comp:betagamma-degree4}
The bar differential $d^4: \bar{B}^4 \to \bar{B}^3$ is determined by
iterated simple-pole coefficient extraction along collision divisors.
Among the $16$ pure tensor types, the key computation is:

\emph{Type $[\beta|\beta|\beta|\gamma]$.}
Consider the element
\begin{equation}\label{eq:bbbg-element}
\xi = \beta(z_1) \otimes \beta(z_2) \otimes \beta(z_3) \otimes
\gamma(z_4) \otimes \eta_{12} \wedge \eta_{13} \wedge \eta_{23}
\wedge \eta_{34}.
\end{equation}
The bar differential acts by summing over all collision divisors $D_{ij}$
($i < j$) and extracting the simple-pole coefficient of the OPE:
\begin{equation}\label{eq:bbbg-differential}
d(\xi) = \sum_{i < j} \operatorname{Coeff}_{D_{ij}}
\bigl[\mu_{ij}(a_i \otimes a_j) \otimes
\operatorname{Res}^{\log}_{D_{ij}}\omega\bigr].
\end{equation}

Since $\beta(z)\beta(w) \sim 0$, the coefficients at $D_{12}$, $D_{13}$,
and $D_{23}$ vanish: these correspond to $\beta$--$\beta$ collisions
with no singular OPE. Only the divisors involving a $\beta$--$\gamma$
collision can contribute. The relevant divisors are $D_{14}$, $D_{24}$,
and $D_{34}$.

\emph{Coefficient at $D_{34}$ (first).}
Set $\epsilon = z_3 - z_4$. The OPE gives
$\beta(z_3)\gamma(z_4) \sim 1/(z_3 - z_4)$, so:
\begin{align*}
\operatorname{Coeff}_{\epsilon^{-1}}
\bigl[\beta(z_3)\gamma(z_4)\bigr]\cdot
\operatorname{Res}^{\log}_{D_{34}}(\eta_{34})
&=1\cdot\mathbb{1} \\
&= \mathbb{1}.
\end{align*}
After this coefficient extraction, the surviving form factor involves
$\eta_{12} \wedge \eta_{13} \wedge \eta_{23}$, restricted to
$\overline{C}_3(X)$ with variables $z_1, z_2, z_3$. But by the
Arnold relation
\begin{equation}\label{eq:arnold-3pt-recall}
\eta_{12} \wedge \eta_{23} + \eta_{23} \wedge \eta_{31}
+ \eta_{31} \wedge \eta_{12} = 0,
\end{equation}
one has $\eta_{12} \wedge \eta_{13} \wedge \eta_{23} = 0$ in
$\Omega^3_{\log}(\overline{C}_3)$ (since $\eta_{13} = -\eta_{31}$
and the three $2$-fold products span a $2$-dimensional space).

Consequently, the coefficient at $D_{34}$ produces a
degree-$3$ element multiplied by a vanishing $3$-form, giving
zero.

\emph{Coefficient at $D_{24}$.}
Setting $\epsilon = z_2 - z_4$, the $\beta(z_2)\gamma(z_4)$ OPE
contributes $\mathbb{1}$. After collapsing $z_2 = z_4$, the surviving
form involves $\eta_{12}, \eta_{13}, \eta_{23}$ with $z_2$ replaced
by $z_4$. By a similar Arnold relation argument, the surviving
$3$-form on three points $(z_1, z_3, z_4)$ equals
$\eta_{14} \wedge \eta_{13}$ (here we use the relation
$\eta_{12}|_{z_2 = z_4} = \eta_{14}$ and
$\eta_{23}|_{z_2 = z_4} = \eta_{43} = -\eta_{34}$). The result is:
\begin{equation}\label{eq:bbbg-D24-residue}
\operatorname{Coeff}_{D_{24}}(\xi) =
-\beta(z_1) \otimes \beta(z_3) \otimes \eta_{14} \wedge \eta_{13}.
\end{equation}
The sign arises from the Koszul convention: removing position~2
from the tensor product $\beta_1 \otimes \beta_2 \otimes \beta_3
\otimes \gamma_4$ and contracting $\beta_2 \gamma_4 \to \mathbb{1}$
requires moving $\beta_2$ past $\beta_3$, which introduces no sign
(both are even/bosonic).

\emph{Coefficient at $D_{14}$.}
By the same mechanism, the $\beta(z_1)\gamma(z_4)$ collision gives:
\begin{equation}\label{eq:bbbg-D14-residue}
\operatorname{Coeff}_{D_{14}}(\xi) =
\beta(z_2) \otimes \beta(z_3) \otimes \eta_{12} \wedge \eta_{23}.
\end{equation}

\emph{Total differential.} Combining:
\begin{equation}\label{eq:bbbg-total}
d(\xi) = \beta(z_2) \otimes \beta(z_3) \otimes \eta_{12}
\wedge \eta_{23}
- \beta(z_1) \otimes \beta(z_3) \otimes \eta_{14}
\wedge \eta_{13}.
\end{equation}
Both terms lie in $\bar{B}^2(\beta\gamma)$ and are exact
(they are boundaries of degree-$3$ elements involving mixed
$\beta\gamma$ types), so $H^3(\bar{B}(\beta\gamma))$ is acyclic in this range.
\end{computation}

\begin{computation}[Nonzero differential types at degree~4]
\label{comp:degree4-nonzero-types}
Among the $16$ pure tensor types at degree~4, one classifies which
differentials are nonzero. The differential $d^4$ acts on type
$[a_1|a_2|a_3|a_4]$ by computing the simple-pole coefficient along
each pair $(i,j)$; the coefficient is nonzero only when
$a_i = \beta$ and $a_j = \gamma$ (or vice versa), giving a nonzero
OPE pole.

\begin{center}
\begin{tabular}{lcc}
\toprule
Type & Nonzero coefficient channels & $d \neq 0$? \\
\midrule
$[\beta|\beta|\beta|\beta]$ & none & No \\
$[\gamma|\gamma|\gamma|\gamma]$ & none & No \\
$[\beta|\beta|\beta|\gamma]$ & $D_{14}, D_{24}, D_{34}$ & Yes \\
$[\gamma|\beta|\beta|\beta]$ & $D_{12}, D_{13}, D_{14}$ & Yes \\
$[\beta|\gamma|\beta|\beta]$ & $D_{12}, D_{23}, D_{24}$ & Yes \\
$[\beta|\beta|\gamma|\beta]$ & $D_{13}, D_{23}, D_{34}$ & Yes \\
$[\gamma|\gamma|\gamma|\beta]$ & $D_{14}, D_{24}, D_{34}$ & Yes \\
$[\beta|\gamma|\gamma|\gamma]$ & $D_{12}, D_{13}, D_{14}$ & Yes \\
$[\gamma|\beta|\gamma|\gamma]$ & $D_{12}, D_{23}, D_{24}$ & Yes \\
$[\gamma|\gamma|\beta|\gamma]$ & $D_{13}, D_{23}, D_{34}$ & Yes \\
$[\beta|\gamma|\beta|\gamma]$ & $D_{12}, D_{14}, D_{23}, D_{34}$ & Yes \\
$[\gamma|\beta|\gamma|\beta]$ & $D_{12}, D_{14}, D_{23}, D_{34}$ & Yes \\
$[\beta|\beta|\gamma|\gamma]$ & $D_{13}, D_{14}, D_{23}, D_{24}$ & Yes \\
$[\gamma|\gamma|\beta|\beta]$ & $D_{13}, D_{14}, D_{23}, D_{24}$ & Yes \\
$[\beta|\gamma|\gamma|\beta]$ & $D_{12}, D_{13}, D_{24}, D_{34}$ & Yes \\
$[\gamma|\beta|\beta|\gamma]$ & $D_{12}, D_{13}, D_{24}, D_{34}$ & Yes \\
\bottomrule
\end{tabular}
\end{center}
A type with $k$ copies of $\beta$ and $4-k$ copies of $\gamma$ has exactly
$k(4-k)$ nonzero coefficient channels, one for each mixed pair. Here $D_{ij}$
denotes the bar differential extracting the OPE between positions $i$
and $j$; for the $\beta\gamma$ system, $D_{ij}$ is nonzero precisely
when the two slots carry complementary fields, since both
$\beta(z)\gamma(w) \sim 1/(z-w)$ and
$\gamma(z)\beta(w) \sim -1/(z-w)$ have simple poles.
Thus the only types with $d = 0$ are the pure $[\beta^4]$ and
$[\gamma^4]$ types. All mixed types have nontrivial bar differential.
\end{computation}

\subsection{Degree~5 and the acyclicity pattern}

At degree~5, the $32$ pure tensor types follow the same pattern:
$[\beta^5]$ and $[\gamma^5]$ have $d = 0$, while all $30$ mixed
types have nontrivial differentials. The derivative types (130 of
them) similarly split into exact and boundary components.

\begin{proposition}[Twisted Koszul acyclicity for
\texorpdfstring{$\beta\gamma$}{beta-gamma};
\ClaimStatusConditional]
\label{prop:betagamma-bar-acyclicity}
\index{beta-gamma system@$\beta\gamma$ system!bar acyclicity}
Let $\tau\colon \bar B(\beta\gamma_\lambda)\to\beta\gamma_\lambda$
be the canonical twisting cochain. The twisted Koszul complex
\[
K(\tau)
\;:=\;
\beta\gamma_\lambda\otimes_\tau\bar B(\beta\gamma_\lambda)
\]
is acyclic off the vacuum:
\[
H^n(K(\tau))=0\quad(n\ge1),\qquad H^0(K(\tau))\cong\mathbb C.
\]
The bare bar coalgebra is not acyclic; its cohomology is the
Koszul-dual coalgebra $(\beta\gamma_\lambda)^{\mathrm i}$.
\end{proposition}

\begin{proof}
Filter $K(\tau)$ by conformal weight. The associated graded complex is
the ordinary Koszul resolution of the polynomial algebra
$\mathrm{Sym}(V)$ with $V=\mathbb C\langle\beta,\gamma\rangle$:
\[
\cdots \longrightarrow
\mathrm{Sym}(V)\otimes \Lambda^r(V^*)
\xrightarrow{\;\iota_{\langle\beta,\gamma\rangle}\;}
\mathrm{Sym}(V)\otimes \Lambda^{r-1}(V^*)
\longrightarrow \cdots .
\]
The contraction is induced by the nondegenerate residue pairing
$\langle\beta,\gamma\rangle=1$, hence the associated graded complex
has cohomology only in the vacuum. The perturbation from the full
chiral bar differential preserves the filtration and has the same
leading residue; the standard filtered perturbation argument gives
the same cohomology for $K(\tau)$.

Removing the left $\beta\gamma$-module factor removes the Koszul
resolution. What remains is the bar coalgebra, whose cohomology is
$(\beta\gamma)^{\mathrm i}$ by
Theorem~\ref{thm:betagamma-bar-cohomology}.\qedhere
\end{proof}

%===================================================================================
% GENUS-1 CURVATURE OF THE BETA-GAMMA SYSTEM
%===================================================================================

\section{Genus-1 curvature of the \texorpdfstring{$\beta\gamma$}{beta-gamma} system}
\label{sec:betagamma-genus1}
\index{beta-gamma system@$\beta\gamma$ system!genus-1 curvature|textbf}
\index{curvature!genus-1 beta-gamma@genus-1 $\beta\gamma$}

We compute the genus-$1$ curvature $m_0^{(1)}(\beta\gamma)$ of the
$\beta\gamma$ system on the elliptic curve
$E_\tau = \mathbb{C}/(\mathbb{Z} + \tau\mathbb{Z})$.

\subsection{The \texorpdfstring{$\beta\gamma$}{beta-gamma} propagator on \texorpdfstring{$E_\tau$}{E tau}}

On the torus $E_\tau$, the genus-$0$ propagator
$P^{(0)}(z,w) = 1/(z-w)$ is replaced by the
elliptic propagator. For the $\beta\gamma$ system with conformal
weights $(\lambda, 1-\lambda)$, the fields are sections of
$K_X^\lambda$ and $K_X^{1-\lambda}$ respectively
(Construction~\ref{const:geometric-beta-gamma}). On $E_\tau$,
$K_{E_\tau} \cong \mathcal{O}_{E_\tau}$, so one needs a flat
line bundle $\mathcal{L}_\zeta$ parameterized by
$\zeta \in \mathrm{Pic}^0(E_\tau)$ to distinguish $\beta$ and
$\gamma$ globally.

\begin{definition}[Elliptic \texorpdfstring{$\beta\gamma$}{beta-gamma} propagator]
\label{def:elliptic-betagamma-propagator}
\index{propagator!elliptic $\beta\gamma$}
The $\beta\gamma$ propagator on $E_\tau$ twisted by the flat
line bundle $\mathcal{L}_\zeta$ is:
\begin{equation}\label{eq:elliptic-betagamma-propagator}
P^{(1)}(z,w|\tau,\zeta) = \frac{\theta_1'(0|\tau) \,
\theta_1(z - w + \zeta|\tau)}
{\theta_1(z-w|\tau)\, \theta_1(\zeta|\tau)}
\end{equation}
where $\theta_1(u|\tau) = -\sum_{n \in \mathbb{Z}}
(-1)^n q^{(n+1/2)^2/2} e^{2\pi i (n+1/2)u}$ is the odd Jacobi
theta function, $q = e^{2\pi i \tau}$, and
$\theta_1'(0|\tau) = \partial_u \theta_1(u|\tau)|_{u=0}$.
\end{definition}

\begin{remark}[Structure of the propagator]
The propagator~\eqref{eq:elliptic-betagamma-propagator} has the
following properties:
\begin{enumerate}
\item A simple pole at $z = w$ with residue~$1$
 (from $\theta_1(0|\tau) = 0$, $\theta_1'(0|\tau) \neq 0$).
\item Quasi-periodicity: $P^{(1)}(z+1,w) = P^{(1)}(z,w)$ and
 $P^{(1)}(z+\tau,w) = e^{-2\pi i \zeta} P^{(1)}(z,w)$.
\item At $\zeta = 0$, the propagator reduces to
 $\theta_1'(0|\tau)/\theta_1(z-w|\tau)$, which is the derivative of
 the prime form on $E_\tau$.
\end{enumerate}
The parameter $\zeta$ corresponds physically to a
Wilson line (background gauge field) for the $U(1)$ symmetry
$\beta \mapsto e^{i\alpha}\beta$,
$\gamma \mapsto e^{-i\alpha}\gamma$.
\end{remark}

\subsection{Curvature computation}

\begin{proposition}[Genus-1 curvature; \ClaimStatusConditional]
\label{prop:betagamma-genus1-curvature}
\index{curvature!genus-1 beta-gamma@genus-1 $\beta\gamma$|textbf}
The genus-$1$ curvature element of the $\beta\gamma$ system with
conformal weight parameter $\lambda$ is:
\begin{equation}\label{eq:betagamma-genus1-curvature}
m_0^{(1)}(\beta\gamma, \lambda)
= \frac{6\lambda^2 - 6\lambda + 1}{12} \cdot E_2(\tau)
= \frac{c_{\beta\gamma}}{24} \cdot E_2(\tau)
\end{equation}
where $c_{\beta\gamma} = 2(6\lambda^2 - 6\lambda + 1)$ is the
central charge (Theorem~\ref{thm:beta-gamma-stress}), and
$E_2(\tau)$ is the \emph{quasi-modular} Eisenstein series of
weight~$2$: not a holomorphic modular form, it transforms as
\[
E_2\!\left(\tfrac{a\tau+b}{c\tau+d}\right)
= (c\tau+d)^2 E_2(\tau) + \tfrac{6c(c\tau+d)}{\pi i}.
\]
\end{proposition}

\begin{proof}
\emph{Step~1: Laurent expansion of the elliptic propagator.}
Expand $P^{(1)}(z,w|\tau,\zeta)$ near $z = w$. Write
$u = z - w$ and use the product formula
$\theta_1(u|\tau) = 2q^{1/8}\sin(\pi u)
\prod_{n=1}^\infty (1-q^n)(1-q^n e^{2\pi i u})(1-q^n e^{-2\pi i u})$.
The logarithmic derivative is:
\begin{equation}\label{eq:logderiv-theta1}
\frac{\theta_1'(u|\tau)}{\theta_1(u|\tau)}
= \pi \cot(\pi u) + 4\pi \sum_{n=1}^\infty
\frac{q^n}{1-q^n} \sin(2\pi n u).
\end{equation}
Expanding near $u = 0$:
\[
\frac{\theta_1'(u|\tau)}{\theta_1(u|\tau)}
= \frac{1}{u} - \frac{\pi^2}{3} E_2(\tau) \cdot u + O(u^3)
\]
where the coefficient of $u$ involves the lattice Eisenstein series
$G_2(\tau) = \frac{\pi^2}{3} E_2(\tau)$
(cf.\ Theorem~\ref{thm:heisenberg-genus-one-complete}).

The full propagator~\eqref{eq:elliptic-betagamma-propagator}
has expansion:
\begin{equation}\label{eq:betagamma-propagator-expansion}
P^{(1)}(z,w|\tau,\zeta) = \frac{1}{u}
+ \frac{\theta_1'(\zeta|\tau)}{\theta_1(\zeta|\tau)}
- \frac{\pi^2}{3} E_2(\tau) \cdot u + O(u^2).
\end{equation}
The constant term $\theta_1'(\zeta|\tau)/\theta_1(\zeta|\tau)$
depends on the twist parameter $\zeta$ and is the genus-$1$ analog
of the background charge.

\emph{Step~2: Curvature from the bar differential.}
The genus-$1$ bar differential on $\bar{B}^2(\beta\gamma)$ has the
form $d^{(1)} = d^{(0)} + \delta^{(1)}$, where $d^{(0)}$ is the
genus-$0$ differential (simple-pole coefficient extraction for the
rational propagator $1/u$)
and $\delta^{(1)}$ encodes the genus-$1$ correction. The
curvature $m_0^{(1)}$ measures the failure of $d^{(1)}$ to
square to zero on the cobar side; equivalently, it is the
obstruction to extending the genus-$0$ bar complex to genus~$1$.
This obstruction is the genus-$1$ component $\theta_1$ of the
universal MC class (Theorem~\ref{thm:explicit-theta}):
for $\beta\gamma$, $\theta_1 = \kappa(\beta\gamma) \cdot \mu
\otimes \lambda_1$, with $\kappa(\beta\gamma) =
6\lambda^2 - 6\lambda + 1$.
The Koszul duality $\beta\gamma \leftrightarrow bc$ reverses the
sign of $\Theta$: $\kappa(\beta\gamma) + \kappa(bc) = 0$.

For the degree-$2$ element $[\beta|\gamma] \otimes \eta_{12}$,
the genus-$1$ differential extracts:
\begin{align}\label{eq:genus1-bar-betagamma-deg2}
d^{(1)}([\beta|\gamma] \otimes \eta_{12}^{(1)})
&=
\operatorname{Coeff}_{u^{-1}}
\bigl[P^{(1)}(u|\tau,\zeta)/u\bigr]
\cdot \mathbb{1}
\end{align}
where $\eta_{12}^{(1)} = d\log\theta_1(z_1 - z_2|\tau)$ is the
elliptic logarithmic form and $u=z_1-z_2$. After the rational
double-pole term is subtracted, the normalized elliptic kernel records
the coefficient of $u^{-1}$ in $P^{(1)}/u$. From
\eqref{eq:betagamma-propagator-expansion},
$P^{(1)}/u = 1/u^2 + (\text{const})/u + \ldots$,
so the genus-$1$ correction is the constant term in the Laurent
expansion of $P^{(1)}$, namely
$\theta_1'(\zeta|\tau)/\theta_1(\zeta|\tau)$.

\emph{Step~3: Stress tensor contribution.}
The curvature $m_0^{(1)}$ arises from the conformal anomaly: the
stress tensor $T^{\beta\gamma}$ has a central charge
$c_{\beta\gamma} = 2(6\lambda^2 - 6\lambda + 1)$, and the genus-$1$
partition function $Z^{(1)} = q^{-c/24} \prod_{n=1}^\infty
(\ldots)$ satisfies:
\begin{equation}\label{eq:dlog-Z}
\frac{\partial}{\partial \tau} \log Z^{(1)}_{\beta\gamma}
= \frac{\pi i \, c_{\beta\gamma}}{12} \cdot E_2(\tau).
\end{equation}
By the general genus-$1$ curvature formula
(cf.\ \S\ref{subsec:heisenberg-genus-one-complete} and
the filtered-curved hierarchy), the curvature element is the
coefficient of $\omega_\tau$ (the normalized holomorphic $1$-form
on $E_\tau$) in the B-cycle monodromy defect:
\[
m_0^{(1)} = \frac{c_{\beta\gamma}}{24} \cdot E_2(\tau)
= \frac{6\lambda^2 - 6\lambda + 1}{12} \cdot E_2(\tau). \qedhere
\]
\end{proof}

\begin{proposition}[Obstruction coefficient; \ClaimStatusConditional]
\label{prop:betagamma-obstruction-coefficient}
\index{obstruction coefficient!beta-gamma@$\beta\gamma$}
The obstruction coefficient of the $\beta\gamma$ system is:
\begin{equation}\label{eq:betagamma-kappa}
\kappa(\beta\gamma) = \frac{c_{\beta\gamma}}{2}
= 6\lambda^2 - 6\lambda + 1.
\end{equation}
This satisfies $\kappa(\beta\gamma) + \kappa(bc) = 0$, where
$\kappa(bc) = c_{bc}/2 = -(6\lambda^2 - 6\lambda + 1)$ is the
obstruction coefficient of the Koszul dual $bc$ system, consistent
with the complementarity theorem
(Theorem~\ref{thm:quantum-complementarity-main}).
\end{proposition}

\begin{proof}
The coefficient $\kappa$ is the genus-$1$ scalar projection of the
universal curvature. At genus~$1$ there is no cross-channel
correction, and the curvature formula of
Proposition~\ref{prop:betagamma-genus1-curvature} gives
$m_0^{(1)}=(c_{\beta\gamma}/24)E_2(\tau)$. Since
$\lambda_1=c_1(\mathbb E)$ is represented on the elliptic curve by
the normalized $E_2(\tau)/12$ term, comparison gives
$\kappa(\beta\gamma) = c_{\beta\gamma}/2$.
For genus $g\ge2$, the identity
$m_0^{(g)}(\cA)=\kappa(\cA)\lambda_g$ is used only in the
uniform-weight or harmonic-decoupling scope of
Theorems~\ref{thm:genus-universality} and
\ref{thm:multi-weight-genus-expansion}; a generic
$\beta\gamma_\lambda$ pair has two complementary weights and is not
made uniform by notation alone.

For the Koszul dual $bc$ system,
$c_{bc} = -c_{\beta\gamma}$
(Proposition~\ref{prop:betagamma-bc-koszul-detailed}),
so $\kappa(\beta\gamma)+\kappa(bc)=0$ on the scalar surface. This is
central-charge complementarity for the Koszul pair; it does not by
itself compute the chiral derived centre
$Z^{\mathrm{der}}_{\mathrm{ch}}(\beta\gamma)$.
\end{proof}

\begin{computation}[Curvature table]
\label{comp:betagamma-curvature-table}
The obstruction coefficient $\kappa(\beta\gamma,\lambda)
= 6\lambda^2 - 6\lambda + 1$ and central charge
$c_{\beta\gamma} = 2\kappa$ for special values of $\lambda$:

\begin{center}
\begin{tabular}{cccc}
\toprule
$\lambda$ & $\kappa(\beta\gamma)$ & $c_{\beta\gamma}$ & Physical system \\
\midrule
$0$ & $1$ & $2$ & Weight-$(0,1)$ bosons \\
$1/2$ & $-1/2$ & $-1$ & Symplectic bosons \\
$1$ & $1$ & $2$ & Weight-$(1,0)$ bosons \\
$2$ & $13$ & $26$ & Quadratic differentials \\
\bottomrule
\end{tabular}
\end{center}

The minimum of $\kappa(\lambda) = 6\lambda^2 - 6\lambda + 1$
occurs at $\lambda = 1/2$, where $\kappa = -1/2$. This is the
unique value at which $c_{\beta\gamma} < 0$. The symmetry
$\kappa(\lambda) = \kappa(1-\lambda)$ reflects the $\beta
\leftrightarrow \gamma$ exchange symmetry.
\end{computation}

\begin{computation}[Genus-1 bar differential at degree 2]
\label{comp:betagamma-genus1-bar-deg2}
The explicit genus-$1$ bar differential on the degree-$2$ element
$[\beta|\gamma]$ on $E_\tau$ proceeds as follows.

At genus~$0$, $d^{(0)}([\beta|\gamma] \otimes \eta_{12}) = \mathbb{1}$
(the vacuum), since the coefficient of $(z_1-z_2)^{-1}$ in
$\beta(z_1)\gamma(z_2)$ is~$1$. The factor
$\eta_{12}=d\log(z_1-z_2)$ records the oriented boundary stratum; it
is not a second pole in the closed collision-residue convention.

At genus~$1$, the propagator acquires an $E_2$ correction.
The elliptic propagator
$P^{(1)}(z,w|\tau) = \theta_1'(z-w|\tau)/\theta_1(z-w|\tau)$
(at $\zeta = 0$) has Laurent expansion:
\[
P^{(1)}(z,w|\tau)
= \frac{1}{z-w} - \frac{\pi^2}{3} E_2(\tau)(z-w) + O((z-w)^3).
\]
The difference from the rational propagator is:
\begin{equation}\label{eq:propagator-difference}
P^{(1)}(z,w|\tau) - P^{(0)}(z,w)
= -\frac{\pi^2}{3} E_2(\tau)(z-w) + O((z-w)^3).
\end{equation}
This correction contributes to the genus-$1$ bar differential at
the next conformal weight level. Specifically, when acting on a
conformal descendant $[\partial\beta|\gamma]$ or $[\beta|\partial\gamma]$,
the $E_2(\tau)$ correction gives a nonzero genus-$1$ contribution:
\[
\delta^{(1)}([\partial\beta|\gamma] \otimes \eta_{12}^{(1)})
\propto E_2(\tau) \cdot \mathbb{1},
\]
confirming that the $E_2(\tau)$ genus-$1$ correction arises from
the difference between the elliptic and rational propagators.
\end{computation}

%===================================================================================
% SPECTRAL SEQUENCE FOR BETA-GAMMA KOSZULNESS
%===================================================================================

\section{Spectral sequence for \texorpdfstring{$\beta\gamma$}{beta-gamma}
Koszulness}
\label{sec:betagamma-spectral-sequence}
\index{spectral sequence!beta-gamma Koszulness@$\beta\gamma$ Koszulness|textbf}

The Koszulness of the $\beta\gamma$ system is the exactness of the
twisted Koszul resolution
$\beta\gamma\otimes_{\tau}\bar B(\beta\gamma)$ together with the
identification
$H^\bullet\bar B(\beta\gamma)=(\beta\gamma)^{\mathrm i}$
(Theorem~\ref{thm:betagamma-bar-cohomology}). Both statements are
controlled by the conformal-weight filtration. The argument parallels
the non-abelian preview of \S\ref{sec:frame-nonabelian-preview}, with
the symplectic pairing $\langle \beta, \gamma \rangle = 1$ replacing
the Killing form.

\subsection{Conformal weight filtration}

\begin{definition}[Conformal weight filtration]
\label{def:conformal-weight-filtration}
\index{filtration!conformal weight}
Define the decreasing filtration $F^p \bar{B}^n(\beta\gamma)$ by
total conformal weight:
\begin{equation}\label{eq:conf-weight-filtration}
F^p \bar{B}^n = \bigl\{
a_1 \otimes \cdots \otimes a_n \otimes \omega \in \bar{B}^n
\;\big|\;
h(a_1) + \cdots + h(a_n) \geq p \bigr\}
\end{equation}
where $h(a_i)$ is the conformal weight of the generator $a_i$
(so $h(\beta) = \lambda$, $h(\gamma) = 1-\lambda$,
$h(\partial^k \beta) = \lambda + k$, etc.).
\end{definition}

The bar differential preserves the filtration: since the OPE
$\beta(z)\gamma(w) \sim 1/(z-w)$ has conformal weight~$0$ on the
right-hand side, coefficient extraction does not increase the total
conformal weight. Therefore $(F^p \bar{B}^*, d)$ is a filtered
coalgebra and the associated spectral sequence converges to the
Koszul-dual coalgebra. After tensoring with $\beta\gamma$ and
twisting by $\tau$, the same filtration gives the Koszul resolution
spectral sequence.

\subsection{\texorpdfstring{$E_1$}{E1} page}

\begin{proposition}[\texorpdfstring{$E_1$}{E1} page; \ClaimStatusProvedHere]
\label{prop:betagamma-E1-page}
For the twisted Koszul resolution, the $E_1$ page of the conformal
weight spectral sequence is:
\begin{equation}\label{eq:betagamma-E1}
E_1^{p,q} =
H^{p+q}\!\bigl(\mathrm{gr}^p(
\beta\gamma\otimes_\tau\bar{B}(\beta\gamma)), \dzero\bigr)
\end{equation}
where $\dzero$ is the leading-order (weight-preserving) part of
the bar differential. The complex
$(\mathrm{gr}^p(\beta\gamma\otimes_\tau\bar B), \dzero)$ is the
classical Koszul complex
\begin{equation}\label{eq:classical-koszul}
\cdots \to \mathrm{Sym}^n(V) \otimes \Lambda^n(V^*)
\xrightarrow{\;\iota\;} \mathrm{Sym}^{n-1}(V)
\otimes \Lambda^{n-1}(V^*) \to \cdots
\end{equation}
where $V = \mathbb{C}\langle \beta, \gamma \rangle$ is the
$2$-dimensional space of generators and $\iota$ is the
contraction map.
\end{proposition}

\begin{proof}
The associated graded $\mathrm{gr}^p \bar{B}$ forgets the
singular part of the OPE and retains only the leading-order
pole structure. At leading order, the $\beta\gamma$ OPE
$\beta(z)\gamma(w) \sim 1/(z-w)$ reduces to the symplectic
pairing $\langle \beta, \gamma \rangle = 1$, and the bar
differential becomes the contraction:
\[
\dzero(a_1 \otimes \cdots \otimes a_n)
= \sum_{i < j} \langle a_i, a_j \rangle \cdot
a_1 \otimes \cdots \widehat{a_i} \cdots \widehat{a_j}
\cdots \otimes a_n.
\]
This is the differential of the Koszul complex
$\mathrm{Sym}(V) \otimes \Lambda(V^*)$ associated to the
identity map $V^* \to V$ (via the symplectic pairing). Since
$\dim V = 2$ and the pairing is nondegenerate, this Koszul
complex is exact in positive degrees (a standard result from
commutative algebra; see~\cite{LV12} \S2.1).
\end{proof}

\begin{computation}[\texorpdfstring{$E_1$}{E_1} page through degree~3]
\label{comp:betagamma-E1-explicit}
\index{spectral sequence!E1 page@$E_1$ page}
Explicitly, at the minimum conformal weight:

\emph{Degree~0.} The vacuum class survives:
$E_1^{0,0}=\mathbb C$.

\emph{Degree~1.} The bar coalgebra has cogenerators represented by
$\beta$ and $\gamma$; after finite duality these become the fermionic
generators $b$ and $c$ of $bc_\lambda$. In the twisted resolution,
the left $\mathrm{Sym}(V)$ factor contracts them by the usual Koszul
differential.

\emph{Degree~2.} The mixed classes are contracted by the coefficient
$\langle\beta,\gamma\rangle=1$. The pure cycles
$[\beta\otimes\beta]$ and $[\gamma\otimes\gamma]$ do not imply
acyclicity of the bare bar coalgebra; they are the next
Koszul-dual coalgebra classes. In the twisted resolution they lie in
the image of the Koszul differential from
\(\mathrm{Sym}(V)\otimes\Lambda^2(V^*)\).

\emph{Degree~3.} Since $\dim V=2$, the leading exterior factor
$\Lambda^3(V^*)$ vanishes in the classical Koszul resolution. Chiral
descendants enter in higher conformal filtration and are controlled
by the same contraction.
\end{computation}

\begin{proposition}[Spectral sequence collapse;
\ClaimStatusConditional]
\label{prop:betagamma-ss-collapse}
\index{spectral sequence!collapse!beta-gamma@$\beta\gamma$}
The conformal weight spectral sequence for the twisted Koszul
resolution of $\beta\gamma_\lambda$ collapses to the vacuum. The
corresponding bare-bar spectral sequence identifies
$H^\bullet\bar B(\beta\gamma_\lambda)$ with the Koszul-dual
coalgebra.
\end{proposition}

\begin{proof}
By Proposition~\ref{prop:betagamma-E1-page}, the twisted-resolution
$E_1$ page is computed by the classical Koszul complex of the
$2$-dimensional symplectic vector space
$(V, \langle\cdot,\cdot\rangle)$ with
$\langle\beta,\gamma\rangle = 1$. This complex is exact in
positive degrees (the pairing is nondegenerate), so
$E_1^{p,q} = 0$ for $p + q \geq 1$ at the leading filtration
level.

The higher filtration levels contribute to $E_1^{p,q}$ only
for $q > 0$ (conformal descendants). The $d_1$ differential
on these terms is induced by the subleading terms in the OPE,
i.e., by the regular part of $\beta(z)\gamma(w)$ beyond the
simple pole. Since the $\beta\gamma$ OPE has
\emph{no regular terms} (the full OPE is exactly
$\beta(z)\gamma(w) = 1/(z-w)$ with no subleading corrections),
the $d_1$ differential on conformal descendants is entirely
determined by the Virasoro action.

The Virasoro module structure on conformal descendants is
semisimple for generic $\lambda$ (the Verma module
$M_{h,c}$ is irreducible for generic $h$ and $c$).
Semisimplicity implies that the higher Koszul differentials
$d_r$ for $r \geq 2$ vanish: the $E_1$ exactness propagates
to $E_2$, and the spectral sequence degenerates.

Concretely, the $E_2$ page satisfies:
\[
E_2^{p,q} = \begin{cases}
\mathbb{C} & (p,q) = (0,0), \\
0 & \text{otherwise},
\end{cases}
\]
and $E_2 = E_\infty$ for the twisted resolution. For the bare bar
coalgebra, the same spectral sequence leaves the diagonal
Koszul-dual coalgebra described in
Theorem~\ref{thm:betagamma-bar-cohomology}.
The collapse holds for all $\lambda$ because the nondegeneracy
of the symplectic pairing $\langle \beta, \gamma \rangle = 1$
is independent of $\lambda$ (the parameter $\lambda$ affects
only the conformal weights, not the leading OPE coefficient).
\end{proof}

%===================================================================================
% SYMPLECTIC BOSON LAMBDA=1/2 SPECIALIZATION
%===================================================================================

\section{Symplectic boson \texorpdfstring{$\lambda = 1/2$}{lambda = 1/2}
specialization}
\label{sec:symplectic-boson-bar}
\index{symplectic boson|textbf}
\index{beta-gamma system@$\beta\gamma$ system!symplectic boson ($\lambda=1/2$)}

At $\lambda = 1/2$, both fields $\beta$ and $\gamma$ have conformal
weight $1/2$, the central charge is $c_{\beta\gamma} = -1$, and the
stress tensor takes the antisymmetric form
$T = \frac{1}{2}(\normord{\beta\partial\gamma}
- \normord{\partial\beta \cdot \gamma})$.
This is the \emph{symplectic boson} system. It is not the $c=-2$
symplectic-fermion vertex algebra, although both occur in logarithmic
conformal field theory.

\subsection{Finite symmetry of the bar
complex}

\begin{definition}[Symplectic rotation and parity]
\label{def:Z2-action-symplectic}
\index{Z2 symmetry@$\mathbb{Z}_2$ symmetry!symplectic boson}
At $\lambda = 1/2$, the two generators have the same conformal
weight and the OPE pairing is skew. The raw exchange
$\beta\leftrightarrow\gamma$ is anti-symplectic. The symplectic
rotation
\[
s(\beta)=\gamma,\qquad s(\gamma)=-\beta
\]
is a vertex-algebra automorphism of order~$4$, and its square
\[
\pi=s^2:\qquad \beta\mapsto-\beta,\qquad \gamma\mapsto-\gamma
\]
is the parity $\mathbb Z_2$ action used below.
\end{definition}

The parity action on the bar complex is
\begin{equation}\label{eq:Z2-on-bar}
\pi(a_1 \otimes \cdots \otimes a_n \otimes \omega)
= (-1)^n a_1 \otimes \cdots \otimes a_n \otimes \omega
\end{equation}
on pure generator tensors of bar degree~$n$. The bar differential
contracts two generators and therefore preserves parity, so
$\bar{B}(\beta\gamma)$ decomposes into parity eigenspaces
$\bar{B} = \bar{B}^+ \oplus \bar{B}^-$.

\subsection{Explicit degree-2 computation at \texorpdfstring{$\lambda = 1/2$}{lambda = 1/2}}

\begin{computation}[Degree-2 bar differential at \texorpdfstring{$\lambda = 1/2$}{lambda = 1/2}]
\label{comp:symplectic-degree2}
At $\lambda = 1/2$, both $\beta$ and $\gamma$ have weight $1/2$,
so the degree-$2$ bar complex has no distinction by conformal
weight between $\beta$ and $\gamma$ insertions. The pure tensor
types and their differentials are:
\begin{align*}
d(\beta \otimes \beta) &= 0, &
d(\gamma \otimes \gamma) &= 0, \\
d(\beta \otimes \gamma) &= +\mathbb{1}, &
d(\gamma \otimes \beta) &= -\mathbb{1}.
\end{align*}
Under parity, all four degree-$2$ tensors are even, while the
degree-$1$ cogenerators are odd and the vacuum is even. The displayed
differential preserves this parity:
\[
d(\bar B^2)^+ \subset \bar B^0{}^+,\qquad
d(\bar B^1)^- \subset 0.
\]
The symplectic rotation $s$ gives a finer order-$4$ symmetry, but the
$\mathbb Z_2$ equivariant bar statement below uses only~$\pi$.
\end{computation}

\subsection{Equivariant bar cohomology}

\begin{proposition}[\texorpdfstring{$\mathbb{Z}_2$}{Z_2}-equivariant bar cohomology;
\ClaimStatusConditional]
\label{prop:symplectic-equivariant-cohomology}
\index{bar cohomology!equivariant}
At $\lambda = 1/2$, the $\mathbb{Z}_2$-equivariant bar
cohomology of the $\beta\gamma$ system is the invariant part of the
Koszul-dual coalgebra:
\begin{equation}\label{eq:equivariant-cohomology}
H^\bullet(\bar{B}(\beta\gamma_{1/2}))^{\mathbb{Z}_2}
\;\cong\;
\bigl((\beta\gamma_{1/2})^{\mathrm i}\bigr)^{\mathbb{Z}_2}.
\end{equation}
The corresponding equivariant twisted Koszul resolution has only the
vacuum cohomology, over characteristic zero and for a semisimple
finite-group action.
\end{proposition}

\begin{proof}
The bar differential commutes with the finite symmetry on the chosen
equivariant model, so the action descends to cohomology. By
Theorem~\ref{thm:betagamma-bar-cohomology}, this cohomology is
$(\beta\gamma_{1/2})^{\mathrm i}$; taking invariants gives
\eqref{eq:equivariant-cohomology}. Over $\mathbb C$ the invariants
functor for a finite group is exact, so applying it to the twisted
Koszul resolution of Proposition~\ref{prop:betagamma-bar-acyclicity}
preserves its vacuum-only cohomology.\qedhere
\end{proof}

\begin{remark}[Symplectic boson versus symplectic fermion]
\label{rem:symplectic-logarithmic}
\index{logarithmic CFT!symplectic boson}
\index{symplectic fermion!not symplectic boson}
The point $\lambda=1/2$ in the bosonic $\beta\gamma$ family is the
rank-one symplectic boson: $c_{\beta\gamma}=-1$ and
$\kappa(\beta\gamma)=-1/2$. Its Verdier/Koszul partner is the
fermionic $bc_{1/2}$ system, the complex-fermion point with
$c_{bc}=1$ and $\kappa(bc)=1/2$. Thus the numerical fixed point of
the weight involution $\lambda\mapsto1-\lambda$ is not a
Koszul self-duality point; the bosonic and fermionic statistics remain
different.

The $c=-2$ triplet algebra $\mathcal W(2)$ belongs to the
symplectic-fermion side
\textup{(}Example~\ref{ex:triplet-logarithmic}\textup{)}, not to the
symplectic-boson parity orbifold. The finite
$\mathbb Z_2$ action on the symplectic-boson bar complex
\textup{(}Computation~\ref{comp:symplectic-degree2}\textup{)} still
produces an equivariant bar category, but this category is not
identified here with the $c=-2$ triplet category. Logarithmic
phenomena therefore enter this chapter only with the following
firewall: the parent $\beta\gamma_{1/2}$ bar complex remains the
bosonic class-$\mathsf C$ free-field complex; triplet
non-semisimplicity is a symplectic-fermion/orbifold phenomenon treated
in the logarithmic-module chapter.
\end{remark}

The $\beta\gamma/bc$ duality is the two-generator chiral lift of
$\mathrm{Com}^! = \mathrm{Lie}$; the equivariant structure at
$\lambda = 1/2$ is a finite-symmetry refinement of the symplectic
boson bar complex, not an identification with the $c=-2$ triplet
algebra.

%% ================================================================
%% Holographic test case
%% ================================================================
\section{The \texorpdfstring{$\beta\gamma$}{beta-gamma} system as holographic test case}
\label{sec:betagamma-holographic}
\index{beta-gamma system@$\beta\gamma$ system!holographic test case}

The free $\beta\gamma$ system furnishes the simplest non-trivial testing
ground for the modular homotopy holography package. Interval
compactification of a 3d $\mathcal{N}=2$ free chiral multiplet
produces the full chiral algebra, and the resulting modular test
package isolates the genus dependence in factorization homology.

\subsection{Interval compactification}

\begin{proposition}[Interval compactification produces the full
 $\beta\gamma$ algebra {\cite{CDG20}, \S4.2}; \ClaimStatusProvedElsewhere]
\label{prop:betagamma-interval-compactification}
\index{interval compactification!beta-gamma@$\beta\gamma$}
Consider the holomorphic twist of a 3d\/ $\mathcal{N}=2$ free chiral multiplet
on $[0,1] \times \Sigma$, equipped with Neumann boundary conditions at both
endpoints. Let $\varphi$ denote the boundary boson and $\psi^{(1)}$ the
first descendant of the fermionic bulk field. Define
\[
 \beta(z) \;:=\; \varphi(z), \qquad
 \gamma(z) \;:=\; \int_0^1 \psi^{(1)}.
\]
Then:
\begin{enumerate}
\item The field $\gamma$ is $Q$-closed: the Stokes boundary term vanishes
 because $\psi$ satisfies Neumann conditions on both boundary components.
\item The OPE $\beta(z)\gamma(0) \sim 1/z$ is inherited from the bulk
 Poisson bracket between $\varphi$ and $\psi$.
\item The interval compactification realizes the full $\beta\gamma$ chiral
 algebra, rather than a single boundary half.
\end{enumerate}
\end{proposition}

\begin{proof}
The holomorphic twist of the 3d $\mathcal{N}=2$ chiral multiplet
$(\varphi, \psi)$ on $[0,1] \times \Sigma$ places $\varphi$ as a scalar
and $\psi$ as a $(0,1)$-form along $\Sigma$ with a first-order kinetic
term along the interval direction. Neumann conditions at $t=0$ and $t=1$
mean that $\psi|_{\partial} = 0$, so the Stokes boundary term in
$Q\gamma = Q\bigl(\int_0^1 \psi^{(1)}\bigr) = \int_0^1 \overline{\partial}_\Sigma \psi^{(1)} + [\psi^{(1)}]_0^1$
vanishes. The singular OPE $\beta(z)\gamma(0) \sim 1/z$ follows from the
bulk propagator $\langle \varphi(z,t)\, \psi(0,t') \rangle$ integrated
over the interval. See \cite{CDG20}, \S4.2 for the complete calculation.
\end{proof}

\subsection{The free modular test package}

\begin{definition}[Free $\beta\gamma$ modular test package]
\label{def:betagamma-modular-test-package}
\index{modular test package!beta-gamma@$\beta\gamma$}
Let $V_{\beta\gamma} = \mathrm{Sym}\bigl(\mathbb{C}[\partial]\beta
\oplus \mathbb{C}[\partial]\gamma\bigr)$ denote the $\beta\gamma$ vertex
algebra with $\lambda$-bracket $\{\beta_\lambda \gamma\} = 1$.
The \emph{free modular test object} is the tuple
\[
 \mathrm{ModEnv}_{\beta\gamma}^{\mathrm{free}}
 \;:=\;
 \bigl(
 V_{\beta\gamma},\;
 V_{\beta\gamma}^{!},\;
 H_{g,n}^{\beta\gamma},\;
 \Theta_{g,n}^{\beta\gamma},\;
 F_{\beta\gamma},\;
 I_{\beta\gamma}^{\mathrm{def}}
 \bigr)
\]
where:
\begin{itemize}
\item $V_{\beta\gamma}^{!} \cong \mathcal{F}$ is the Koszul dual free
 fermion (Theorem~\ref{thm:cobar-fermions});
\item $H_{g,n}^{\beta\gamma}$ denotes the genus-$g$, $n$-pointed
 factorization homology (chiral homology) of $V_{\beta\gamma}$;
\item $\Theta_{g,n}^{\beta\gamma}$ is the tautological action of the
 mapping class group on $H_{g,n}^{\beta\gamma}$;
\item $F_{\beta\gamma}$ is the factorization structure on $\mathrm{Ran}(X)$;
\item $I_{\beta\gamma}^{\mathrm{def}}$ collects all higher-genus deformation
 data.
\end{itemize}
In this package the higher \emph{local OPE} $A_\infty$ operations are
\emph{strict} (all local $m_k = 0$ for $k \geq 3$). Transferred
modular contact classes, such as the class-$\mathsf C$ quartic
coefficient, belong to the modular shadow calculus rather than to
new singular terms in the vertex-algebra OPE.
\end{definition}

\subsection{The determinant line}

\begin{definition}[$\beta\gamma$ determinant line]
\label{def:betagamma-determinant-line}
\index{determinant line!beta-gamma@$\beta\gamma$}
Fix a conformal weight parameter $\lambda$ and an auxiliary line bundle
$L$ on a smooth projective curve $\Sigma_g$ of genus~$g$. The
\emph{$\beta\gamma$ determinant line} is
\begin{equation}\label{eq:betagamma-det-line}
 \mathcal{D}^{\beta\gamma}_{g,\lambda}(L)
 \;:=\;
 \bigl(\!\det R\Gamma(\Sigma_g,\, K_{\Sigma_g}^{\lambda} \otimes L)\bigr)^{-1}
 \;\otimes\;
 \bigl(\!\det R\Gamma(\Sigma_g,\, K_{\Sigma_g}^{1-\lambda} \otimes L^\vee)\bigr)^{-1}.
\end{equation}
Here $K_{\Sigma_g}$ is the canonical bundle and $\det R\Gamma$ denotes the
determinant of cohomology in the sense of Knudsen--Mumford.
\end{definition}

\begin{remark}[Bosonization and the determinant line]
\label{rem:betagamma-determinant-bosonization}
At $\lambda = 1/2$ and $L = \mathcal{O}$, the Riemann--Roch theorem gives
$\dim H^0(\Sigma_g, K^{1/2}) = \dim H^1(\Sigma_g, K^{1/2})$ for
$g \geq 1$, so the two tensor factors in
$\mathcal{D}^{\beta\gamma}_{g,1/2}(\mathcal{O})$ are \emph{isomorphic}
(but not canonically so; the isomorphism depends on a choice of theta
characteristic). This reflects the bosonization ambiguity of the
symplectic boson at $\lambda = 1/2$.
\end{remark}

\begin{proposition}[Mumford exponent complementarity]
\label{prop:mumford-exponent-complementarity}
\ClaimStatusProvedHere
\index{Mumford isomorphism!exponent complementarity}
\index{determinant line!Mumford exponent}
The Mumford exponents $e(n) := 6n^2 - 6n + 1$ satisfy:
\begin{equation}\label{eq:mumford-exponent-complementarity}
 e(\lambda) + e(1 - \lambda)
 \;=\;
 c_{\beta\gamma}(\lambda)
 \;=\;
 2(6\lambda^2 - 6\lambda + 1).
\end{equation}
At $\lambda = 1$ \textup{(}standard weights\textup{)},
$e(1) + e(0) = 1 + 1 = 2 = c_{\beta\gamma}(1)$. The total
Mumford exponent of the determinant line
$\mathcal{D}^{\beta\gamma}_{g,\lambda}$ is $-c_{\beta\gamma}$, and
complementarity at the Mumford level gives
$\det(\beta\gamma) \otimes \det(bc) \simeq \mathcal{O}$: the
exponents cancel because $c_{\beta\gamma} + c_{bc} = 0$.
\end{proposition}

\begin{proof}
\leavevmode
\begin{enumerate}[label=\textup{(\roman*)}]
\item \emph{Exponent identity.}
Expand directly:
\[
 e(\lambda) = 6\lambda^2 - 6\lambda + 1,
 \qquad
 e(1-\lambda) = 6(1-\lambda)^2 - 6(1-\lambda) + 1
 = 6\lambda^2 - 6\lambda + 1.
\]
Hence $e(\lambda) + e(1-\lambda) = 2(6\lambda^2 - 6\lambda + 1)$.
The central charge of the $\beta\gamma$ system at conformal weight
$\lambda$ is $c_{\beta\gamma}(\lambda) = 2(6\lambda^2 - 6\lambda + 1)$
\textup{(}two fields of spins $\lambda$ and $1-\lambda$\textup{)},
establishing~\eqref{eq:mumford-exponent-complementarity}.

\item \emph{Standard weights.}
At $\lambda = 1$: $e(1) = 6 - 6 + 1 = 1$ and $e(0) = 1$, so
$e(1) + e(0) = 2 = c_{\beta\gamma}(1)$.

\item \emph{Determinant-line cancellation.}
The Mumford isomorphism gives
$\det R\Gamma(\Sigma_g, K^n) \simeq \lambda_1^{e(n)}$
where $\lambda_1 = \det R\Gamma(\Sigma_g, K)$ is the Hodge
determinant. The $\beta\gamma$ determinant line at weight $\lambda$
has total exponent $-(e(\lambda) + e(1-\lambda)) = -c_{\beta\gamma}(\lambda)$.
For the $bc$ system at the \emph{same} weight parameter, the central
charge is $c_{bc}(\lambda) = -c_{\beta\gamma}(\lambda)$
\textup{(}fermionic sign\textup{)}. Therefore
\[
 \det(\beta\gamma)_\lambda \otimes \det(bc)_\lambda
 \;\simeq\;
 \lambda_1^{-c_{\beta\gamma} + c_{\beta\gamma}}
 \;=\;
 \lambda_1^{0}
 \;\simeq\;
 \mathcal{O},
\]
which is the Mumford-level manifestation of boson--fermion
complementarity.
\qedhere
\end{enumerate}
\end{proof}

\subsection{Modular state spaces and the determinant conjecture}

\begin{conjecture}[Free-field realization of modular state spaces]
\label{conj:betagamma-modular-state-spaces}
\index{modular state spaces!beta-gamma conjecture@$\beta\gamma$ conjecture}
\ClaimStatusConjectured
For the free $\beta\gamma$ theory at conformal weight $\lambda$, the
genus-$g$ partition function is conjecturally the regularized determinant
\begin{equation}\label{eq:betagamma-partition-determinant}
Z_g^{\beta\gamma}(\lambda)
= \det\nolimits' R\Gamma(\Sigma_g, K^{\lambda})^{-1}
 \cdot \det\nolimits' R\Gamma(\Sigma_g, K^{1-\lambda})^{-1},
\end{equation}
where $\det'$ denotes the zeta-regularized determinant of the Dolbeault
Laplacian and $K = K_{\Sigma_g}$ is the canonical bundle. This is a
section of the determinant line
$\mathfrak{D}^{\beta\gamma}_{g,\lambda}$ of
Definition~\textup{\ref{def:betagamma-determinant-line}}.

In the modular envelope framework, this conjecture amounts to: the
modular state space $\mathcal{H}_g^{\beta\gamma}$ is one-dimensional,
spanned by $Z_g^{\beta\gamma}$, and the factorization-homology integral
reduces to the Quillen determinant. Concretely:
\begin{enumerate}
\item The genus-$g$ partition function $Z_g^{\beta\gamma}$ is a canonical
 section of the determinant line
 $\mathcal{D}^{\beta\gamma}_{g,\lambda}(\mathcal{O})$ over $\mathcal{M}_g$,
 obtained by integrating the Dolbeault complexes
 $\bigl(\Omega^{0,\bullet}(\Sigma_g, K^{\lambda}),\, \overline{\partial}\bigr)$
 and
 $\bigl(\Omega^{0,\bullet}(\Sigma_g, K^{1-\lambda}),\, \overline{\partial}\bigr)$
 underlying $\beta$ and~$\gamma$ respectively.
\item The local OPE is free (all higher local operations vanish), but the
 global all-genera content is encoded by the variation of
 $\mathcal{D}^{\beta\gamma}_{g,\lambda}(L)$ over the moduli of curves
 and auxiliary bundles.
\end{enumerate}
\end{conjecture}

\begin{remark}[The $\beta\gamma$ system as a minimal modular test case]
\label{rem:betagamma-cleanest-test}
The $\beta\gamma$ system is a minimal test case for the modular envelope because:
(1)~the OPE is free (single simple pole, no higher singularities);
(2)~all higher local OPE operations vanish
($m_k^{\mathrm{loc}} = 0$ for $k \ge 3$);
(3)~genus dependence enters entirely through determinant-line geometry on curves.
These three properties isolate the global modular contribution from local higher-operation corrections.
\end{remark}


%% ================================================================
\begin{remark}[A minimal holographic test case]
\label{rem:betagamma-holographic-atom}
\index{beta-gamma system@$\beta\gamma$ system!minimal holographic test case}
\index{holomorphic-topological!beta-gamma test case@$\beta\gamma$ test case}
Local OPE triviality ($m_k^{\mathrm{loc}} = 0$ for $k \ge 3$) forces
the analytic genus dependence into the global geometry of curves. The
$\beta\gamma$ system is therefore a minimal example in which the
boundary algebra and the Koszul-dual defect sector are explicit, while
the universal closed sector is the separate chiral Hochschild object
$Z^{\mathrm{der}}_{\mathrm{ch}}(\beta\gamma)$, not the Koszul dual
algebra $bc_\lambda$.
\end{remark}

\section{From 3d \texorpdfstring{$\mathcal{N}=2$}{N=2} to the full chiral algebra}

The interval-compactification construction of
Proposition~\ref{prop:betagamma-interval-compactification} identifies the
$\beta\gamma$ system as a boundary algebra. The resulting fields give
the boundary input for a holographic datum; the intrinsic closed-sector slot is
the chiral derived center, not the Koszul dual algebra.

\begin{definition}[Boundary $\beta\gamma$ holographic package]
\label{def:betagamma-holographic-package}
\index{holographic package!beta-gamma@$\beta\gamma$}
The \emph{boundary $\beta\gamma$ holographic package} consists of the
following data:
\begin{enumerate}
\item The vertex algebra
 $V_{\beta\gamma} = \mathrm{Sym}\bigl(\mathbb{C}[\partial]\beta
 \oplus \mathbb{C}[\partial]\gamma\bigr)$ with
 $\lambda$-bracket $\{\beta_\lambda \gamma\} = 1$.
\item The boundary fields
 $\beta(z) := \varphi(z)$ and
 $\gamma(z) := \int_0^1 \psi^{(1)}$,
 where $\varphi$ is the boundary boson and $\psi^{(1)}$ the first
 descendant of the bulk fermion (Proposition~\ref{prop:betagamma-interval-compactification}).
\item The $Q$-closure of $\gamma$, which holds by Stokes' theorem:
 the bulk fermion $\psi$ vanishes on both Neumann boundary
 components, so $Q\gamma = [\psi^{(1)}]_0^1 = 0$.
\item The singular OPE $\beta(z)\gamma(0) \sim 1/z$, inherited from
 the bulk Poisson bracket between $\varphi$ and $\psi$.
\end{enumerate}
As a boundary modular test object, the $\beta\gamma$ system has
$\mu_{0,2}^{\beta\gamma} \neq 0$ (the genus-$0$ two-point pairing is
non-degenerate) but $\mu_{g,n}^{\beta\gamma} = 0$ for all higher local
OPE operations ($m_k^{\mathrm{loc}} = 0$ for $k \geq 3$). The
analytic genus dependence therefore enters through factorization
homology on the curve:
\begin{equation}\label{eq:betagamma-state-spaces}
 H_{g,n}^{\beta\gamma}(L_1, \ldots, L_n)
 \;=\;
 \int_{\Sigma_{g,n}} \mathrm{Fact}_{\beta\gamma}[L_1, \ldots, L_n].
\end{equation}
\end{definition}

\subsection{The determinant line}

For the free Gaussian model, determinant-line control is the natural
candidate for the factorization-homology answer. The determinant line
$\mathcal{D}_{g,\lambda}^{\beta\gamma}(L)$
(Definition~\ref{def:betagamma-determinant-line})
is the line on which the conjectural genus-$g$ state lives.

\begin{conjecture}[Determinantal control of genus-$g$ state spaces;
\ClaimStatusConjectured]
\label{conj:betagamma-determinant-control}
\index{determinant line!genus control}
For each $g \geq 0$, the genus-$g$ state space $H_g^{\beta\gamma}$
of the free $\beta\gamma$ system is governed by the determinant
line~\eqref{eq:betagamma-det-line}: the factorization homology
$\int_{\Sigma_g} \mathrm{Fact}_{\beta\gamma}$ is quasi-isomorphic to
$\mathcal{D}_{g,\lambda}^{\beta\gamma}$
as a graded line over the moduli space $\mathcal{M}_g$.
\end{conjecture}

\begin{remark}[Why determinantal]
\label{rem:betagamma-why-determinantal}
This is the expected output of the Gaussian path integral: for a free
field theory, the partition function is a ratio of functional
determinants, which on a compact Riemann surface are finite-dimensional
determinant lines. The two factors in~\eqref{eq:betagamma-det-line}
correspond to the $\beta$-sector ($K^\lambda \otimes L$) and the
$\gamma$-sector ($K^{1-\lambda} \otimes L^\vee$), reflecting the
Serre-dual pairing between the two fields.
\end{remark}

\subsection{The formal genus series}

\begin{definition}[Formal $\beta\gamma$ genus series]
\label{def:betagamma-holographic-partition-function}
\index{partition function!beta-gamma holographic@$\beta\gamma$ holographic}
The \emph{formal $\beta\gamma$ genus series} is the power series in
$\hbar$ defined by
\begin{equation}\label{eq:betagamma-partition-function}
 Z_{\beta\gamma}(\hbar)
 \;:=\;
 \sum_{g \geq 0} \hbar^{2g-2}\,
 \chi_{\mathrm{gr}}\bigl(H_g^{\beta\gamma}\bigr),
\end{equation}
where $\chi_{\mathrm{gr}}$ denotes the graded Euler characteristic.
\end{definition}

\begin{remark}[Determinantal genus expansion for $\beta\gamma$]
\label{rem:betagamma-holographic-atom-principle}
\index{beta-gamma system@$\beta\gamma$ system!determinantal genus expansion}
Local OPE triviality ($m_k^{\mathrm{loc}} = 0$ for
$k \geq 3$) strips away local combinatorial corrections, leaving the genus expansion
$Z_{\beta\gamma}(\hbar)$ governed by the global geometry of
moduli. Every term $\chi_{\mathrm{gr}}(H_g^{\beta\gamma})$ is, by
Conjecture~\ref{conj:betagamma-determinant-control}, a determinant-line
invariant on $\mathcal{M}_g$, so the generating function
\eqref{eq:betagamma-partition-function} records how
factorization homology interacts with stable-curve compactification,
with no local nonlinear OPE correction terms. Any failure of the modular
package at the $\beta\gamma$ level is therefore a failure of the
global gluing axioms, not of local algebra.

Compare with the Heisenberg system (Overture), where
$c_{\mathrm{Heis}} = 1$ already detects modular curvature. The
Heisenberg is \emph{not} Koszul self-dual
($\cH_k^! = \mathrm{Sym}^{\mathrm{ch}}(V^*)$, not~$\cH_k$),
but its abelian OPE structure makes the two sides difficult to
distinguish at the level of shadow invariants. The $\beta\gamma$
system, also not self-dual ($(\beta\gamma)^{!} \simeq bc$),
separates the boundary algebra, the Koszul-dual defect sector, and
the derived-center closed-sector slot.
\end{remark}


%% ================================================================
\section{The \texorpdfstring{$\beta\gamma$}{beta-gamma} quartic birth and rank-one rigidity}
\label{sec:betagamma-quartic-birth}
\index{beta-gamma system@$\beta\gamma$ system!quartic birth}
\index{nonlinear shadow!quartic archetype}

The preceding sections established that the $\beta\gamma$ system is
locally free: all higher local $A_\infty$ operations vanish at the
level of the strict vertex-algebra OPE. The class-$\mathsf C$ datum is
not a local OPE correction. It is the transferred modular contact
datum of the standard conformal-weight family
$\beta\gamma_\lambda$, with $S_3=0$, $S_4=-5/12$, and no tail beyond
degree~$4$. The internal weight-changing line is the rank-one
degeneration where the shadow tower vanishes.
Its depth decomposition $d = 1 + d_{\mathrm{arith}} + d_{\mathrm{alg}}
= 1 + 1 + 2 = 4$
(Theorem~\ref{thm:depth-decomposition}): the $\mathrm{U}(1)$
Epstein zeta $\varepsilon^1_s = 4\zeta(2s)$ gives
$d_{\mathrm{arith}} = 1$ (one critical line), while the
standard-family contact datum gives $d_{\mathrm{alg}} = 2$
(Chapter~\ref{chap:arithmetic-shadows}).

Appendix~\ref{subsec:nms-betagamma} contains the full shadow calculus
for the three primitive archetypes. The $\beta\gamma$ construction in
this chapter is the class-$\mathsf C$ instance.

\subsection{Tower construction for \texorpdfstring{$\beta\gamma$}{beta-gamma}}
\label{subsec:betagamma-tower-construction}
\index{shadow obstruction tower!beta-gamma construction}

\paragraph{Step~1: Degree~$2$.}
On the weight-changing line, the genus-$0$ curvature vanishes ($m_0 = 0$) and the
degree-$2$ shadow is trivial:
$\Theta_{\beta\gamma}^{\leq 2}\big|_{\text{w.c.}} = 0$.
(The global modular characteristic $\kappa = c/2 = 6\lambda^2 - 6\lambda + 1$
is generically nonzero; the vanishing here is specific to the
$\lambda$-marginal slice.)

\paragraph{Step~2: Degree~$3$, cubic shadow vanishes.}
On the weight-changing line, the Maurer--Cartan equation is linear
and no cubic vertex exists:
$\mathfrak{C}_{\beta\gamma} = 0$ and $o_3 = 0$.
$\Theta_{\beta\gamma}^{\leq 3}|_{\mathrm{w.c.}} = 0$.

\paragraph{Step~3: Degree~$4$, standard-family contact.}
The class-$\mathsf C$ standard family has quartic coefficient
$S_4=-5/12$. On the internal weight-changing line the possible
quartic scalar is
$\mu_{\beta\gamma}=\langle\eta,m_3(\eta,\eta,\eta)\rangle$, and
rank-one abelian rigidity gives
$\mu_{\beta\gamma}=0$
(Corollary~\ref{cor:nms-betagamma-mu-vanishing}).

\paragraph{Step~4: Termination.}
For the standard conformal-weight family, $S_r=0$ for $r\ge5$.
On the weight-changing line the stronger statement holds:
$\Theta_{\beta\gamma}^{\le r}|_{\mathrm{w.c.}}=0$ for all $r\ge2$.

\paragraph{Step~5: Arithmetic face.}
In the reduced completion convention,
$W(\beta\gamma) = \{1,1\}$, so
$\operatorname{ek}(\beta\gamma) = 0$ and
$S_{\beta\gamma}(u) = 2\,\zeta(u)\,\zeta(u{+}1)$.
The standard conformal-weight family is \emph{exact Euler--Koszul}
despite having shadow depth~$4$. In contrast, the direct
sum $\cH \oplus \cH$ shares the same weight multiset
$\{1,1\}$, the same character $\eta(q)^{-2}$, and the same
$\operatorname{ek} = 0$, but has shadow depth~$2$
(Gaussian, not contact). This is the decisive witness that
$\operatorname{ek}$ depends on the character while
$\kappa_d$ depends on the OPE
(Corollary~\ref{cor:two-faces-theta}).

\medskip
The construction is formalized in
Theorem~\ref{thm:betagamma-quartic-birth} and
Corollary~\ref{cor:betagamma-postnikov-termination} below.

\subsection{The weight/contact slice}

\begin{definition}[Weight/contact slice; \ClaimStatusDefinitional]
\label{def:betagamma-weight-contact-slice}
\index{weight/contact slice!beta-gamma@$\beta\gamma$}
Let $L_{\mathrm{w.c.}}\subset V_{\beta\gamma}$ be the
one-dimensional weight-changing line spanned by the class $\eta$ that
shifts conformal weight $\lambda\mapsto\lambda+\epsilon$. Let
$L_C$ denote the standard conformal-weight class-$\mathsf C$
contact direction, with coordinate~$u$. The
\emph{weight/contact slice} is the smallest cyclic deformation slice
containing $L_{\mathrm{w.c.}}$ and $L_C$.
\end{definition}

\subsection{Quartic birth}

\begin{theorem}[\texorpdfstring{$\beta\gamma$}{beta-gamma} quartic birth]
\label{thm:betagamma-quartic-birth}
\ClaimStatusProvedHere
\index{quartic birth!beta-gamma@$\beta\gamma$}
On the weight/contact slice
\textup{(}Definition~\textup{\ref{def:betagamma-weight-contact-slice}}\textup{)}
the cubic shadow vanishes on the standard class-$\mathsf C$ direction
and the quartic functional is the cyclic transfer of the ternary
operation:
\[
 \mathfrak{C}_{\beta\gamma}\big|_{L_C} = 0,
 \qquad
 \mathfrak{Q}_{\beta\gamma} = \operatorname{cyc}(m_3).
\]
On $L_C=\mathbb C u$,
\[
 \mathfrak{Q}_{\beta\gamma}\big|_{L_C}
 =
 -\frac{5}{12}\,u^4,
 \qquad
 S_2 = 6\lambda^2 - 6\lambda + 1,\quad
 S_3=0,\quad S_4=-\frac{5}{12},\quad S_r=0\ (r\ge5).
\]
On the rank-one weight-changing line,
\[
 \mathfrak{Q}_{\beta\gamma}(\eta,\eta,\eta,\eta)
 \;=\;
 \mu_{\beta\gamma}
 \;:=\;
 \langle \eta,\, m_3(\eta,\eta,\eta) \rangle.
\]
This scalar is zero: $\mu_{\beta\gamma}=0$.
\end{theorem}

\begin{proof}
The mixed simple-pole OPE gives the bar differential but no scalar
cubic coefficient in the class-$\mathsf C$ standard family:
$S_3=0$. In a cyclic Hamiltonian an $m_3$ operation contributes to a
quartic functional, so the first possible contact term is
$\operatorname{cyc}(m_3)$. The restriction to the standard contact
direction $L_C$ is the class-$\mathsf C$ clause of
Theorem~\ref{thm:nms-betagamma-quartic-birth}: the quartic coefficient
is $-5/12$ and the tail vanishes for $r\ge5$. Restricting the same
quartic functional to $L_{\mathrm{w.c.}}=\mathbb C\eta$ gives
$\mu_{\beta\gamma}$; Corollary~\ref{cor:betagamma-mu-vanishing}
gives $\mu_{\beta\gamma}=0$.\qedhere
\end{proof}

\begin{corollary}[$\beta\gamma$ weight-changing line is shadow-trivial]
\label{cor:betagamma-postnikov-termination}
\ClaimStatusProvedHere
\index{shadow obstruction tower!beta-gamma termination@$\beta\gamma$ termination}
On the weight-changing line, the shadow obstruction tower
\textup{(}Definition~\textup{\ref{def:shadow-postnikov-tower})}
of $\beta\gamma$ vanishes in every degree:
\[
 \Theta_{\beta\gamma}^{\le r}\big|_{\mathrm{w.c.}} = 0
 \qquad (r \ge 2).
\]
The class-$\mathsf{C}$ witness is instead the standard
conformal-weight family $\beta\gamma_\lambda$, for which
\[
 S_2 = 6\lambda^2 - 6\lambda + 1,\qquad
 S_3 = 0,\qquad
 S_4 = -\tfrac{5}{12},\qquad
 S_r = 0 \ \ (r \ge 5).
\]
\end{corollary}

\begin{proof}
Rank-one abelian rigidity gives
$\Theta_{\beta\gamma}^{\le r}\big|_{\mathrm{w.c.}} = 0$
for all $r \ge 2$ on the one-dimensional weight-changing slice.
The standard conformal-weight family formulas are the class-$\mathsf C$
part of Theorem~\ref{thm:betagamma-quartic-birth}.
\end{proof}

\begin{proposition}[Class-\texorpdfstring{$\mathsf C$}{C} stratum separation]
\label{prop:betagamma-class-c-stratum-separation}
\ClaimStatusConditional
\index{shadow depth!class C stratum separation}
The equality $S_4=-5/12$ for the standard family
$\beta\gamma_\lambda$ does not assert that an isolated primary line
with $\kappa\neq0$ and $S_4\neq0$ has finite tail. The finite
class-$\mathsf C$ depth is a two-stratum statement: the charged
standard-family contact direction carries the nonzero quartic
coefficient and satisfies $S_r=0$ for $r\ge5$, while the internal
weight-changing line has
\[
 \Theta_{\beta\gamma}^{\le r}\big|_{\mathrm{w.c.}}=0
 \qquad (r\ge2).
\]
\end{proposition}

\begin{proof}
Proposition~\ref{prop:depth-gap-trichotomy} separates the
single-primary-line Riccati regime from the global contact witness.
On a primary line with $\kappa\neq0$, a nonzero quartic coefficient
would make the shadow metric non-square and force an infinite tail.
The $\beta\gamma$ class-$\mathsf C$ witness is not that situation:
the coefficient $S_4=-5/12$ is attached to the standard two-field
conformal-weight family, and the rank-one internal weight-changing
line has zero quartic projection by rank-one abelian rigidity. Thus
the statements $S_4=-5/12$, $S_r=0$ for $r\ge5$, and
$\mu_{\beta\gamma}=0$ are compatible, because they live on different
strata of the weight/contact slice.\qedhere
\end{proof}

\subsection{The \texorpdfstring{$\beta\gamma$}{beta-gamma} primitive kernel}
\label{subsec:betagamma-primitive-kernel}
\index{primitive kernel!beta-gamma@$\beta\gamma$}

\begin{proposition}[\texorpdfstring{$\beta\gamma$}{beta-gamma}
primitive-kernel contact decomposition]
\label{prop:betagamma-primitive-kernel}
\ClaimStatusConditional
For the standard conformal-weight family, the primitive logarithmic
modular kernel
\textup{(}Definition~\textup{\ref{def:primitive-log-modular-kernel})}
of the $\beta\gamma$ system separates the ambient contact corolla from
its rank-one weight-line scalar:
\begin{equation}\label{eq:betagamma-primitive-kernel}
\mathfrak{K}_{\beta\gamma}
\;=\;
K_{0,2}^{\beta\gamma} + K_{0,4}^{\beta\gamma}
+ K_{1,1}^{\beta\gamma} + R_{\mathrm{pf}}^{\beta\gamma},
\end{equation}
where:
\begin{enumerate}[label=\textup{(\roman*)}]
\item $K_{0,2}^{\beta\gamma}$
 is the degree-$2$ corolla from the bilinear $\beta\gamma$ pairing;
\item $K_{0,3}^{\beta\gamma} = 0$
 \textup{(}cubic vanishes on the weight-changing line\textup{)};
\item $K_{0,4}^{\beta\gamma}
 = \operatorname{cyc}(m_3)$
 is the quartic contact corolla from the transferred operation~$m_3$,
 with $\mu_{\beta\gamma}
 := \langle\eta, m_3(\eta,\eta,\eta)\rangle = 0$
 \textup{(}Corollary~\textup{\ref{cor:nms-betagamma-mu-vanishing})};
\item $K_{1,1}^{\beta\gamma}$ is the genus-$1$ corolla;
\item $R_{\mathrm{pf}}^{\beta\gamma}$ is the rigid-cutting residue
 from the contact stratum of Mok's log-FM boundary
 \textup{(}Theorem~\textup{\ref{thm:logfm-modular-cocomposition})}.
\end{enumerate}
\end{proposition}

\begin{proof}
The weight-changing OPE is linear in $\lambda$ and produces
no cubic vertex: $\ell_2^{\mathrm{tr}} = 0$ on the weight-changing
line, so $K_{0,3} = 0$. The first possible transferred contact
operation is $m_3$ (from the composite
$:$$\beta\partial\gamma$$:$ interaction), giving the ambient corolla
$K_{0,4} = \operatorname{cyc}(m_3)$. Its scalar projection on the
weight-changing line is zero. Higher operations
$m_n$ for $n \geq 4$ vanish on this slice by conformal-weight
counting: the weight-changing class $\eta$ has weight~$2$, and
$m_n(\eta^{\otimes n})$ would require a weight-$(2n-2)$ output
in a weight-$2$ space, which is impossible for $n \geq 4$
by the polynomial filtration on the OPE. Therefore $K_{0,n} = 0$
for $n \geq 5$ on the weight-changing branch. The rigid-cutting
residue $R_{\mathrm{pf}}$ belongs to the ambient contact package:
it is present on the standard class-$\mathsf C$ surface with
$S_4=-5/12$ and has zero scalar projection on the rank-one
weight-changing line.
\end{proof}

\begin{proposition}[\texorpdfstring{$\beta\gamma$}{beta-gamma}
primitive shell equations under log-FM contact compatibility]
\label{prop:betagamma-primitive-shell}
\ClaimStatusConditional
\index{shell equation!beta-gamma@$\beta\gamma$}
Assume the log-FM rigid-cutting boundary contributes to the standard
class-$\mathsf C$ contact channel with coefficient $S_4=-5/12$.
Then the genus-$2$ shell equation for $\beta\gamma$ exhibits all
three channels of the primitive master equation:
\begin{equation}\label{eq:betagamma-shell-genus2}
R_2(\mathfrak{K}_{\beta\gamma})
\;=\;
\underbrace{\Delta_{\mathrm{ns}}(K_{1,1}^{\beta\gamma})}_{\text{BV}}
\;+\;
\underbrace{\tfrac{1}{2}\,
[K_{1,1}^{\beta\gamma},\,K_{1,1}^{\beta\gamma}]}_{\text{gauge-trivial}}
\;+\;
\underbrace{\sum_{\rho \in \operatorname{Rig}_2}
R_\rho^{\beta\gamma}(K_{0,2}^{\beta\gamma})}_{\text{rigid cutting}}.
\end{equation}
The separating bracket vanishes by antisymmetry on the rank-one
weight-changing branch. The BV channel is nonzero; the rigid-cutting
channel is the conditional ambient contact contribution of the
standard family.
\end{proposition}

\begin{proof}
The BV self-contraction $\Delta_{\mathrm{ns}}(K_{1,1})$ is nonzero
(it produces the standard genus-raising term). The separating
bracket $[K_{1,1}, K_{1,1}] = 0$ by rank-$1$ antisymmetry on the
weight-changing line. The rigid-cutting channel is the ambient
class-$\mathsf C$ contribution: the quartic contact corolla
$K_{0,4}$ composes with the degree-$2$ corolla $K_{0,2}$ through the
planted-forest boundary of $\overline{M}_{2,n}$. This is the channel absent
from both Heisenberg
(Proposition~\ref{prop:heisenberg-primitive-shell})
and affine
(Proposition~\ref{prop:affine-primitive-shell}), and its
presence is the signature of the standard-family contact class.
\end{proof}

\begin{computation}[\texorpdfstring{$\beta\gamma$}{beta-gamma}
branch BV action]
\label{comp:betagamma-branch-bv}
\index{branch BV action!beta-gamma@$\beta\gamma$}
On the weight-changing line the branch-active quotient
\textup{(}Definition~\textup{\ref{def:reduced-branch-master-action})}
is one-dimensional: $V_{\beta\gamma}^{\mathrm{br}} = \mathbb{C}\,\eta$
with coordinate~$x$. Since $\mu_{\beta\gamma} = 0$
(Corollary~\ref{cor:nms-betagamma-mu-vanishing}), the branch
master action reduces to
\begin{equation}\label{eq:betagamma-branch-action}
S_{\beta\gamma}^{\mathrm{br}}(x)
\;=\;
\tfrac{1}{2}\,\kappa(\beta\gamma)\,x^2.
\end{equation}
The quartic vertex $\frac{1}{4!}\mu_{\beta\gamma} x^4 = 0$
is killed by rank-one abelian rigidity.
Despite the ambient quartic corolla of the standard family, the
contact invariant on the weight-changing line projects to zero in the
branch BV packet. The spectral discriminant and metaplectic half-density
coincide with the Heisenberg pattern:
\begin{equation}\label{eq:betagamma-metaplectic}
\Delta_{\beta\gamma}(x)\big|_{\mathrm{w.c.}} = 1 - \kappa(\beta\gamma)\,x,
\qquad
\delta_{\beta\gamma}(x)\big|_{\mathrm{w.c.}} = (1 - \kappa(\beta\gamma)\,x)^{1/2}.
\end{equation}
The branch packet on the weight-changing line is Gaussian, but the
full $\beta\gamma$ primitive kernel is \emph{not} Gaussian: the
quartic corolla $K_{0,4}$ and the rigid-cutting residue
$R_{\mathrm{pf}}$ carry non-scalar data visible on other slices.
\end{computation}

\subsection{Explicit computations on the weight-changing line}
\label{subsec:betagamma-explicit-weight-line}

The computations below give the rank-one weight-line restriction of
Theorem~\ref{thm:betagamma-quartic-birth}. They isolate the
$\beta\gamma$ OPE, the homotopy transfer formula, and the
complementarity potential on the line
$L_{\mathrm{w.c.}}=\mathbb C\eta$.

\subsubsection{The weight-changing class \texorpdfstring{$\eta$}{eta}}
\label{sssec:betagamma-eta-class}
\index{weight-changing class!beta-gamma@$\beta\gamma$}

The $\beta\gamma$ system has OPE
$\beta(z)\gamma(w) \sim 1/(z-w)$ with conformal weights
$(h_\beta, h_\gamma) = (\lambda, 1-\lambda)$. The stress tensor
(Theorem~\ref{thm:beta-gamma-stress}) is
\[
 T_\lambda(z)
 \;=\;
 (1-\lambda)\,\normord{\beta(z)\,\partial\gamma(z)}
 \;-\;
 \lambda\,\normord{\partial\beta(z)\,\gamma(z)},
\]
depending linearly on $\lambda$. The weight-changing deformation
shifts $\lambda \mapsto \lambda + \epsilon$, producing a one-parameter
family $T_{\lambda+\epsilon}$ of stress tensors.

\begin{definition}[The marginal class $\eta$]
\label{def:betagamma-eta-marginal}
The \emph{weight-changing class}
$\eta \in H^1(\operatorname{Def}_{\mathrm{cyc}}(\beta\gamma))$
is the tangent vector to this family:
\[
 \eta
 \;=\;
 \frac{\partial T_{\lambda+\epsilon}}{\partial \epsilon}
 \bigg|_{\epsilon=0}
 \;=\;
 \normord{\beta\,\partial\gamma}
 \;-\;
 \normord{(\partial\beta)\,\gamma}.
\]
This is a weight-$(2,0)$ field; it is marginal (conformal dimension~$2$)
and lies in the cyclic deformation complex since both $\beta$ and
$\gamma$ participate symmetrically.
\end{definition}

The cyclic pairing on $\eta$ is controlled by the central charge.
The $TT$ OPE gives
\[
 T_\lambda(z)\,T_\mu(w)
 \;\sim\;
 \frac{c(\lambda,\mu)/2}{(z-w)^4}
 \;+\; \cdots
\]
and the central charge of the $\beta\gamma$ system is
$c_{\beta\gamma}(\lambda) = 2(6\lambda^2 - 6\lambda + 1)$.
The quartic pole of $T_\lambda(z)\,T_\lambda(w)$ gives
$c_{\beta\gamma}/2$. Differentiating twice:
\begin{equation}\label{eq:betagamma-eta-pairing}
 \langle \eta, \eta \rangle
 \;=\;
 \frac{\partial^2}{\partial\lambda^2}
 \biggl(\frac{c_{\beta\gamma}(\lambda)}{2}\biggr)
 \;=\;
 \frac{1}{2}\cdot 24
 \;=\; 12,
\end{equation}
since $c_{\beta\gamma}''(\lambda) = 24$ identically. This is the
Hessian $H_{\beta\gamma}$ of the complementarity potential restricted
to the weight-changing line.

\subsubsection{Vanishing of the transferred binary bracket}
\label{sssec:betagamma-m2-vanishing}
\index{transferred bracket!vanishing on weight line}

\begin{lemma}[$\ell_2^{\mathrm{tr}}(\eta,\eta) = 0$;
\ClaimStatusProvedHere]
\label{lem:betagamma-ell2-vanishing}
The transferred binary bracket vanishes on the weight-changing line:
$\ell_2^{\mathrm{tr}}(\eta,\eta) = 0$.
\end{lemma}

\begin{proof}
The deformation $\lambda \mapsto \lambda + \epsilon$ enters
\emph{linearly} in the stress tensor: $T_{\lambda+\epsilon} =
T_\lambda + \epsilon\,\eta$. The OPE
$T_{\lambda+\epsilon}(z)\,T_{\lambda+\epsilon}(w)$ therefore
has coefficients that are at most quadratic in~$\epsilon$.
The \emph{simple pole} (which encodes the Lie bracket)
is $T_{\lambda+\epsilon}(w)/(z-w)$, linear in
$T_{\lambda+\epsilon}$ and hence in $\epsilon$. There is no
$\epsilon^2$ term in the simple pole. Since the transferred binary
bracket $\ell_2^{\mathrm{tr}}$ is extracted from the simple-pole
coefficient, one obtains $\ell_2^{\mathrm{tr}}(\eta,\eta) = 0$.

Equivalently, on a one-dimensional deformation space the
Maurer--Cartan equation
$m_1(\eta) + \tfrac{1}{2}m_2(\eta,\eta) = 0$ is automatically
linear: $m_1$ is the linearized deformation equation (which is the
\emph{only} term), so $m_2(\eta,\eta)$ must vanish identically.
\end{proof}

\subsubsection{The \texorpdfstring{$A_\infty$}{A-infinity} relation
forcing \texorpdfstring{$m_3 = 0$}{m3 = 0}}
\label{sssec:betagamma-m3-vanishing}
\index{homotopy transfer!forcing vanishing of m3@forcing vanishing of $m_3$}

The homotopy transfer formula expresses the transferred ternary
operation $\ell_3^{\mathrm{tr}}$ in terms of lower operations and the
contracting homotopy $h$:
\begin{equation}\label{eq:betagamma-ell3-transfer}
 \ell_3^{\mathrm{tr}}(x,y,z)
 \;=\;
 \ell_2^{\mathrm{tr}}\bigl(h \cdot \ell_2^{\mathrm{tr}}(x,y),\, z\bigr)
 \;+\;
 \ell_2^{\mathrm{tr}}\bigl(x,\, h \cdot \ell_2^{\mathrm{tr}}(y,z)\bigr)
 \;+\; (\text{terms involving } \ell_1).
\end{equation}
On $V_{\beta\gamma} = \mathbb{C}\langle \eta \rangle$ (one-dimensional),
every argument is proportional to~$\eta$. Evaluating on the
triple $(\eta, \eta, \eta)$:

\begin{proposition}[$\ell_3^{\mathrm{tr}}(\eta,\eta,\eta) = 0$;
\ClaimStatusProvedHere]
\label{prop:betagamma-ell3-vanishing}
\index{ternary operation!vanishing on weight line}
The transferred ternary operation vanishes on the weight-changing line:
\[
 \ell_3^{\mathrm{tr}}(\eta,\eta,\eta) \;=\; 0.
\]
\end{proposition}

\begin{proof}
From the displayed homotopy transfer formula, every summand in
$\ell_3^{\mathrm{tr}}(\eta,\eta,\eta)$ factors through
$\ell_2^{\mathrm{tr}}(\eta,\eta)$:
\begin{align*}
 \ell_3^{\mathrm{tr}}(\eta,\eta,\eta)
 &\;=\;
 \ell_2^{\mathrm{tr}}\bigl(
 h \cdot \underbrace{\ell_2^{\mathrm{tr}}(\eta,\eta)}_{=\,0},\,
 \eta\bigr)
 \;+\;
 \ell_2^{\mathrm{tr}}\bigl(
 \eta,\,
 h \cdot \underbrace{\ell_2^{\mathrm{tr}}(\eta,\eta)}_{=\,0}\bigr)
 \;+\; 0 \\[4pt]
 &\;=\; 0 + 0 + 0 \;=\; 0.
\end{align*}
The first two terms vanish because
$\ell_2^{\mathrm{tr}}(\eta,\eta) = 0$
(Lemma~\ref{lem:betagamma-ell2-vanishing}).
The remaining terms from the transfer formula involve either
$\ell_1$-images (which vanish on the cohomology representative~$\eta$)
or higher iterations that again factor through
$\ell_2^{\mathrm{tr}}(\eta,\eta) = 0$.

Hence the quartic contact invariant vanishes:
\[
 \mu_{\beta\gamma}
 \;=\;
 \langle \eta,\, \ell_3^{\mathrm{tr}}(\eta,\eta,\eta) \rangle
 \;=\;
 \langle \eta,\, 0 \rangle
 \;=\; 0.
\]
This is the \emph{homotopy transfer proof} of the vanishing.
The rank-one rigidity proof
(Theorem~\ref{thm:betagamma-rank-one-rigidity} below)
gives the same result by a different mechanism:
on a one-dimensional slice with vanishing brackets,
\emph{all} higher jets are zero.

In the secondary Borcherds language
(Definition~\ref{def:secondary-borcherds}), the vanishing
$\ell_3^{\mathrm{tr}}(\eta,\eta,\eta) = 0$ means that the
$\beta\gamma$~system has \emph{trivial quartic secondary Borcherds
operations}: $j'_{(p,q,r)} = 0$ for all $(p,q,r)$ on the
weight-changing line. The standard conformal-weight family is
nevertheless class~$\mathsf C$: its scalar cubic coefficient is
$S_3=0$, its quartic contact coefficient is $S_4=-5/12$, and its
shadow obstruction tower terminates at weight~$4$.
\end{proof}

\subsubsection{The complementarity potential on the weight line}
\label{sssec:betagamma-potential-weight-line}
\index{complementarity potential!beta-gamma weight line@$\beta\gamma$ weight line}

Assembling the preceding computations, the complementarity potential
restricted to the weight-changing direction takes the form
\begin{equation}\label{eq:betagamma-potential-weight-line}
 \mathfrak{S}_{\beta\gamma}(x)
 \;=\;
 \tfrac{1}{2}\,\langle \eta,\eta \rangle\, x^2
 \;+\;
 \underbrace{\tfrac{1}{3!}\,\langle \eta,\ell_2^{\mathrm{tr}}
 (\eta,\eta)\rangle}_{=\,0}\, x^3
 \;+\;
 \underbrace{\tfrac{1}{4!}\,\langle \eta,\ell_3^{\mathrm{tr}}
 (\eta,\eta,\eta)\rangle}_{=\,0}\, x^4
 \;+\; \cdots
 \;=\;
 \tfrac{1}{2}\,H_{\beta\gamma}\, x^2,
\end{equation}
where $H_{\beta\gamma} = \langle\eta,\eta\rangle = 12$ is the
Hessian~\eqref{eq:betagamma-eta-pairing}. The potential is
\emph{exactly quadratic} on the weight-changing line:
all cubic, quartic, and higher terms vanish.

This means that complementarity is formally fake on this line
(Proposition~\ref{prop:fake-complementarity-criterion}).
The quartic contact invariant
$\mu_{\beta\gamma} = \langle \eta, m_3(\eta^3) \rangle = 0$
is the explicit verification.

\begin{remark}[Off the weight-changing line]
\label{rem:betagamma-off-weight-line}
The vanishing of $\mu_{\beta\gamma}$ is specific to the
\emph{rank-one weight direction}.
Off the weight-changing line, on the full weight/contact
slice that includes transverse deformations, the quartic
shadow $\mathfrak{Q}_{\beta\gamma}$ has standard-family coefficient
$S_4=-5/12$.
The extended contact structure of the $\beta\gamma$ system
(on a two-or-more-dimensional deformation slice) lies beyond
the rank-one abelian rigidity mechanism and requires the full
nonlinear shadow calculus of
Appendix~\ref{subsec:nms-betagamma}.
\end{remark}

\subsection{Vanishing of the quartic contact invariant}

\begin{corollary}[Vanishing of the quartic contact invariant]
\label{cor:betagamma-mu-vanishing}
\ClaimStatusProvedHere
\index{quartic contact invariant!vanishing}
On the weight-changing line of the $\beta\gamma$ system,
\[
 \mu_{\beta\gamma}
 \;:=\;
 \langle \eta,\, m_3(\eta,\eta,\eta) \rangle
 \;=\; 0.
\]
The quartic contact invariant vanishes identically on the weight-changing
deformation.
\end{corollary}

\begin{proof}
First, homotopy transfer gives the vanishing directly.
The transferred binary bracket satisfies $m_2(\eta,\eta) = 0$ on the
one-dimensional weight-changing subspace (the Maurer--Cartan equation
there is linear). Hence
$m_3(\eta,\eta,\eta) = 0$ by the $A_\infty$ relation
\[
 m_3 = m_2 \circ (h \otimes \mathrm{id}) \circ m_2 + \cdots
\]
applied to the vanishing binary input: every summand contains a factor
of $m_2(\eta,\eta) = 0$.

\medskip
Second, the weight-changing line is a one-dimensional cyclic subspace
with vanishing higher brackets; Theorem~\ref{thm:betagamma-rank-one-rigidity}
then kills every higher jet on this line.\qedhere
\end{proof}

\subsection{Rank-one abelian rigidity}

\begin{theorem}[Rank-one abelian rigidity; \ClaimStatusProvedHere]
\label{thm:betagamma-rank-one-rigidity}
\index{rank-one abelian rigidity}
Let $L \subset V_{\cA}$ be a one-dimensional cyclic subspace on which
all transferred higher brackets vanish:
\[
 \ell_n^{\mathrm{tr}}\big|_L = 0
 \qquad (n \ge 2).
\]
Then the restriction of the complementarity potential to $L$ is exactly
quadratic. In particular, complementarity on $L$ is formally fake.
\end{theorem}

\begin{proof}
On such a line the cyclic action expansion collapses to its quadratic
Hessian part. Every higher jet
$\mathrm{Sh}_r(\Theta_{\cA})|_L$ for $r \ge 3$ involves at least one
factor of some $\ell_n^{\mathrm{tr}}$ with $n \ge 2$, and all such
factors vanish by hypothesis. The complementarity potential therefore
reduces to
\[
 \mathfrak{S}_{\cA}\big|_L
 \;=\;
 \tfrac{1}{2}\,\mathrm{Sh}_2(\Theta_{\cA})\big|_L,
\]
which is exactly quadratic.\qedhere
\end{proof}

\subsection{The contact boundary law}

\begin{corollary}[Pure contact boundary law; \ClaimStatusProvedHere]
\label{cor:betagamma-pure-contact-boundary}
\index{contact boundary law!beta-gamma@$\beta\gamma$}
On the weight/contact slice of the $\beta\gamma$ system,
\[
 \xi^*\mathfrak{Q}_{\beta\gamma}
 \;=\;
 H_{\beta\gamma} \star_{P_{\beta\gamma}} \mathfrak{Q}_{\beta\gamma}
 \;+\;
 \mathfrak{Q}_{\beta\gamma} \star_{P_{\beta\gamma}} H_{\beta\gamma}.
\]
There is no cubic tree contribution.
\end{corollary}

\begin{proof}
Apply the quartic boundary recursion
of Corollary~\ref{cor:nms-betagamma-boundary-law} with
$\mathfrak{C}_{\beta\gamma} = 0$: the cubic--quartic cross terms
$\mathfrak{C} \star \mathfrak{Q}$ that appear in the general recursion
are absent, leaving only the quadratic--quartic propagation displayed
above.\qedhere
\end{proof}

\subsection{The trichotomy of primitive archetypes}

\begin{remark}[Uniform explanation of linear examples]
\label{rem:betagamma-uniform-linearity}
Rank-one abelian rigidity
(Theorem~\ref{thm:betagamma-rank-one-rigidity}) explains uniformly why
the Heisenberg level line and the $\beta\gamma$ weight-shift line are
both formally linear: in each case the deformation lies on a
one-dimensional sector with vanishing higher brackets. Real nonlinear
complementarity requires either a non-abelian bracket (as for affine
$\widehat{\mathfrak{sl}}_2$) or a higher-dimensional deformation space
(as for $\beta\gamma$ off the weight-shift line).
\end{remark}

\begin{remark}[Phase index and the contact/quartic archetype]
\label{rem:betagamma-phase-index}
\index{nonlinear phase index!beta-gamma@$\beta\gamma$}
The nonlinear phase index
(Definition~\ref{def:nms-nonlinear-phase-index}) of the standard
conformal-weight family $\beta\gamma_\lambda$ is
$\nu_{\mathrm{nl}}(\beta\gamma) = 4$, the contact/quartic archetype.
The scalar cubic coefficient vanishes and the first nonzero
standard-family contact coefficient is quartic. The internal
weight-changing line is the formally quadratic degeneration of this
family.
Together with Heisenberg ($\nu_{\mathrm{nl}} = \infty$, purely
Gaussian) and affine $\widehat{\mathfrak{sl}}_2$
($\nu_{\mathrm{nl}} = 3$, Lie/tree type), this completes the
\emph{trichotomy of primitive archetypes}
(Theorem~\ref{thm:nms-archetype-trichotomy}).
\end{remark}

\begin{remark}[The $\beta\gamma$ position in the shadow obstruction tower landscape]
\label{rem:betagamma-landscape-position}
\index{shadow obstruction tower!beta-gamma position@$\beta\gamma$ position}
The trichotomy of primitive archetypes is the shadow obstruction tower classification
restricted to the free-field and first-interacting families. In the
full landscape (Table~\ref{tab:shadow-tower-census}), the four
archetype classes are:
\[
\underbrace{\text{G (Gaussian)}}_{\substack{r_{\max}=2\\
\text{Heis, fermion}}} \;\subset\;
\underbrace{\text{L (Lie/tree)}}_{\substack{r_{\max}=3\\
\widehat{\fg}_k}} \;\subset\;
\underbrace{\text{C (contact)}}_{\substack{r_{\max}=4\\
\beta\gamma,\; bc}} \;\subset\;
\underbrace{\text{M (mixed)}}_{\substack{r_{\max}=\infty\\
\text{Vir, }\mathcal{W}_N}}
\]
The $\beta\gamma$ system occupies the unique intermediate position:
the shadow obstruction tower is nontrivial beyond the scalar level but still
terminates at finite degree. This makes it the natural testing ground
for the nonlinear shadow calculus, complex enough to exhibit genuine
higher structure, simple enough for closed-form computation.
\end{remark}


\begin{proposition}[Why $\beta\gamma$ is class~$\mathsf{C}$: standard conformal-weight family]
\label{prop:betagamma-sugawara-class-c}
\ClaimStatusProvedHere
\index{beta-gamma system@$\beta\gamma$ system!class C mechanism}
\index{Sugawara construction!class C depth}
The standard conformal-weight family $\beta\gamma_\lambda$ is class~$\mathsf{C}$
$($shadow depth $r_{\max} = 4)$, with
\[
 S_2 = 6\lambda^2 - 6\lambda + 1,\qquad
 S_3 = 0,\qquad
 S_4 = -\tfrac{5}{12},\qquad
 S_r = 0 \ \ (r \ge 5).
\]
By contrast, the internal weight-changing line is shadow-trivial:
\[
 \Theta_{\beta\gamma}^{\le r}\big|_{\mathrm{w.c.}} = 0
 \qquad (r \ge 2).
\]
\end{proposition}

\begin{proof}
The standard conformal-weight family formulas are
Theorem~\ref{thm:betagamma-quartic-birth} restricted to the
class-$\mathsf C$ contact direction. The weight-changing slice
statement is Corollary~\ref{cor:betagamma-postnikov-termination}.
\end{proof}

\begin{remark}[Completion kinematics of $\beta\gamma$]
\label{rem:betagamma-completion-kinematics}
\index{completion kinematics!beta-gamma@$\beta\gamma$}
The $\beta\gamma$ system has vacuum character
$\chi_{\beta\gamma}(q) = \prod_{n \geq 0}(1-q^{n+1})^{-1}
\cdot \prod_{n \geq 1}(1-q^n)^{-1}$
(two towers, one per generator). The primitive generating series is
$G_{\beta\gamma}(t) = t/(1-t)$: one primitive per reduced weight,
arising from the derivative towers $\partial^r\beta$ and
$\partial^r\gamma$. The Koszul radius is $\rho_K = 1/2$, the
completion entropy is $h_K(\beta\gamma) = \log 2 \approx 0.693$,
and the primitive defect $\Delta_{\beta\gamma}(t) = 0$.

The reduced-weight-$q$ window $K_q(\beta\gamma)$ is a
finite-dimensional complex of dimension $\sim 2^q$ (the ordered-word
species on a one-letter-per-weight alphabet). The bar differential
on $K_q$ is controlled by the contact data: at $q = 4$, the quartic
shadow $\mathfrak{Q}_{\beta\gamma}$ acts as the differential on the
$8$-dimensional complex $K_4$. The vanishing
$\mu_{\beta\gamma} = 0$ is the rank-one weight-line projection of
this window; the standard-family class-$\mathsf C$ datum is still
$S_4=-5/12$ and the tower terminates at degree~$4$.

Kinematically, $\beta\gamma$ has completion entropy
$h_K=\log 2$, while the Virasoro value recorded in the census is
$0.567$. Shadow depth (dynamic complexity) and completion entropy
(kinematic complexity) are independent invariants.
\end{remark}


%% ================================================================
\section{Vortex lines and the explicit line \texorpdfstring{$r$}{r}-matrix}
\label{sec:betagamma-vortex-lines}
\index{vortex line!beta-gamma@$\beta\gamma$}
\index{r-matrix@$r$-matrix!beta-gamma@$\beta\gamma$}

The collision-residue $r$-matrix of the closed chiral algebra is
$r^{\mathrm{coll}}_{\beta\gamma}(z)=0$, because the simple generator
OPE is absorbed by the $d\log$ residue. The $bc$ ghost system has the
same closed convention,
$r^{\mathrm{coll}}_{bc}(z)=0$; its surviving first-order datum is the
graded Clifford ordered contact tensor
$b\otimes c+c\otimes b$. Vortex lines carry a different object: a
rational line $r$-matrix valued in the open line algebra. The two are
not identified.

\subsection{Vortex line modules}

\begin{definition}[Free chiral vortex lines]
\label{def:betagamma-vortex-lines}
\index{vortex line!classification}
For $N \in \mathbb{Z}$, the \emph{vortex line} $V_N$ is the
$\beta\gamma$ line operator characterised by the singular behaviour
\[
 X(z) \;\sim\; z^N \times (\text{regular}),
 \qquad
 \psi(z) \;\sim\; z^{-N} \times (\text{regular}).
\]
As modules over the open line algebra $A_{\mathrm{line}}$:
\begin{enumerate}
\item $V_0 \cong \mathbb{C}$ (the trivial/vacuum line);
\item $V_{N}$ for $N > 0$: $V_N \cong \mathbb{C}[\psi_0,\ldots,\psi_{N-1}]$
 (a polynomial algebra on the $N$ zero-modes of $\psi$);
\item $V_{N}$ for $N < 0$: $V_N \cong \mathbb{C}[X_0,\ldots,X_{|N|-1}]$
 (a polynomial algebra on the $|N|$ zero-modes of $X$).
\end{enumerate}
\end{definition}

\subsection{The explicit \texorpdfstring{$r$}{r}-matrix}

\begin{definition}[The $\beta\gamma$ line $r$-matrix; \ClaimStatusDefinitional]
\label{def:betagamma-r-matrix}
\index{r-matrix@$r$-matrix!explicit formula}
The \emph{rational line $r$-matrix} of the $\beta\gamma$ system is
\begin{equation}\label{eq:betagamma-r-matrix}
 r^{\mathrm{line}}_{\beta\gamma}(z)
 \;=\;
 \sum_{n,m \ge 0}
 (-1)^n \binom{n+m}{m}
 \frac{\psi_n \otimes X_m \;-\; X_n \otimes \psi_m}{z^{n+m+1}}.
\end{equation}
This is a formal distribution in $z^{-1}$ valued in
$A_{\mathrm{line}}\otimes A_{\mathrm{line}}$, where
$A_{\mathrm{line}}=\mathbb{C}[X_n,\psi_n]_{n\ge0}$ is the polynomial
open-line algebra. It is not the closed chiral Koszul dual algebra
$A^!_{\mathrm{ch}}\simeq bc_\lambda$.
\end{definition}

The line $r$-matrix intertwines with translation as follows.

\begin{proposition}[Translation and coproduct]
\label{prop:betagamma-translation-coproduct}
\ClaimStatusProvedHere
\index{translation operator!beta-gamma@$\beta\gamma$}
The translation operator $\tau_z$ acts by
\[
 \tau_z(X_n)
 \;=\;
 \sum_{m=0}^{n} \binom{n}{m} z^m\, X_{n-m},
 \qquad
 \tau_z(\psi_n)
 \;=\;
 \sum_{m=0}^{n} \binom{n}{m} z^m\, \psi_{n-m}.
\]
The coproduct is cocommutative (reflecting the freeness of the
$\beta\gamma$ system):
\[
 \Delta_z(X_n) = \tau_z(X_n) \otimes 1 + 1 \otimes X_n,
 \qquad
 \Delta_z(\psi_n) = \tau_z(\psi_n) \otimes 1 + 1 \otimes \psi_n.
\]
\end{proposition}

\begin{proof}
The translation formula follows from the Taylor expansion of
$X(w+z) = \sum_n X_n (w+z)^{-n-1}$, and similarly for $\psi$.
Cocommutativity of $\Delta_z$ is immediate from the abelian OPE:
$X(z)\psi(w) \sim 1/(z-w)$ has no nonlinear terms, so the Hopf
structure is that of a polynomial algebra (i.e.\ primitively
generated).\qedhere
\end{proof}

\subsection{Fusion of vortex lines}

\begin{example}[Fusion $V_1 \otimes_z V_{-1}$]
\label{ex:betagamma-v1-vm1-ope}
\index{vortex line!fusion}
The simplest nontrivial fusion is $V_1 \otimes_z V_{-1}$.
As a vector space, $V_1 \otimes V_{-1} \cong
\mathbb{C}[\psi_0] \otimes \mathbb{C}[X_0]
\cong \mathbb{C}[\psi_0, X_0]$.
The line $r$-matrix~\eqref{eq:betagamma-r-matrix} specialises to
\[
 r^{\mathrm{line}}_{1,-1}(z)
 \;=\;
 \frac{1}{z}\,\psi_0\, X_0,
\]
which acts as the differential on the tensor product complex. For
$z \neq 0$, the cohomology is
\[
 H^*(V_1 \otimes_z V_{-1})
 \;\cong\;
 \mathbb{C}
 \;\cong\;
 V_0.
\]
That is, the vortex line $V_1$ fuses with $V_{-1}$ to give the
trivial line: line operators form a group under fusion, with
$V_N^{-1} = V_{-N}$.
\end{example}

\begin{remark}[The $\beta\gamma$ avatar of the dg-shifted Yangian]
\label{rem:betagamma-dg-yangian}
\index{dg-shifted Yangian!beta-gamma avatar@$\beta\gamma$ avatar}
The open line algebra
$A_{\mathrm{line}} = \mathbb{C}[X_n,\psi_n]_{n \ge 0}$, equipped with
the coproduct $\Delta_z$ and line
$r$-matrix~\eqref{eq:betagamma-r-matrix}, is the $\beta\gamma$
avatar of the dg-shifted Yangian of Vol~II
(Theorem~\ref{thm:yangian-koszul-dual}). Because the $\beta\gamma$
system is free, the Yangian structure is abelian: the line
$r$-matrix satisfies the classical Yang--Baxter equation with zero
Schouten bracket, and the coproduct is strictly cocommutative.
The vortex lines $\{V_N\}_{N \in \mathbb{Z}}$ form the
\emph{weight lattice} of this abelian Yangian, analogous to the
Verma modules of a non-abelian Yangian.
\end{remark}

\begin{proposition}[$\bar{B}(\cA)$-comodule structure on vortex lines]
\label{prop:betagamma-vortex-comodule}
\ClaimStatusProvedHere
\index{vortex line!bar-comodule structure}
\index{twisting cochain!universal!beta-gamma@$\beta\gamma$}
Each vortex line $V_N$ carries a natural
$\bar{B}^{\mathrm{ch}}(\cA_{\beta\gamma})$-comodule structure: the
coaction
\[
 \rho \;:\; V_N \;\longrightarrow\;
 \bar{B}^{\mathrm{ch}}(\cA_{\beta\gamma}) \otimes V_N
\]
is determined by the line-operator coupling to the boundary chiral
algebra. The rational line $r$-matrix
\eqref{eq:betagamma-r-matrix} is the evaluation of the
\emph{line-evaluation twisting cochain}
$\tau_{\mathrm{line}} \colon \bar{B}^{\mathrm{ch}}(\cA_{\beta\gamma})
\to A_{\mathrm{line}}$ on the spectral parameter~$z$:
\begin{equation}\label{eq:betagamma-twisting-cochain}
 r^{\mathrm{line}}_{\beta\gamma}(z) \;=\; \tau_{\mathrm{line},z}
 \;\in\;
 A_{\mathrm{line}} \otimes A_{\mathrm{line}}\,[\![z^{-1}]\!].
\end{equation}
In particular, the classical Yang--Baxter equation for
$r^{\mathrm{line}}_{\beta\gamma}(z)$ is a
formal consequence of the Maurer--Cartan equation for~$\tau_{\mathrm{line}}$ in the
strict convolution dg~Lie algebra
$\Convstr(\bar{B}^{\mathrm{ch}}(\cA),\, A_{\mathrm{line}})$.
\end{proposition}

\begin{proof}
The boundary action of $\cA_{\beta\gamma}$ on $V_N$ is a chiral
module structure. By the universal property of the reduced chiral
bar construction, the corresponding action map is equivalent to a
twisting cochain
$\tau_{\mathrm{line}}\colon \bar{B}^{\mathrm{ch}}(\cA_{\beta\gamma})
\to A_{\mathrm{line}}$ satisfying the Maurer--Cartan equation in
$\Convstr(\bar{B}^{\mathrm{ch}}(\cA_{\beta\gamma}),A_{\mathrm{line}})$.
Adjunction converts this twisting cochain into the displayed
conilpotent coaction on every vortex module~$V_N$. Evaluating
$\tau_{\mathrm{line}}$ on the two free generators and translating by
$z$ gives exactly the kernel
\eqref{eq:betagamma-r-matrix}. Since
$d\tau_{\mathrm{line}}+\frac12[\tau_{\mathrm{line}},\tau_{\mathrm{line}}]=0$,
the classical Yang--Baxter equation for
$r^{\mathrm{line}}_{\beta\gamma}(z)$ is its two-point component.
\end{proof}

\begin{remark}[Five projections of $\Theta_{\beta\gamma}$]
\label{rem:betagamma-five-identifications}
\index{modular bar--Hamiltonian!beta-gamma projections@$\beta\gamma$ projections}
\index{Chriss--Ginzburg structure principle!beta-gamma@$\beta\gamma$}
The master MC element
$\Theta_{\beta\gamma} \in
\operatorname{MC}(\mathfrak{g}^{\mathrm{amb}}_{\beta\gamma})$
(Definition~\ref{def:modular-bar-hamiltonian}) encodes the boundary
modular bar package of the $\beta\gamma$ system. The five main
theorems below are projections of this object; the physical bulk, when
an open/closed comparison is chosen, is the separate chiral derived
center
$Z^{\mathrm{der}}_{\mathrm{ch}}(\beta\gamma)$
(Remark~\ref{rem:five-from-one}):
\begin{center}
\small
\renewcommand{\arraystretch}{1.3}
\begin{tabular}{lll}
\toprule
\emph{Theorem} & \emph{Projection of $\Theta_{\beta\gamma}$}
 & \emph{Physical content} \\
\midrule
A (bar-cobar) &
 $\Theta_{\beta\gamma}\big|_{\hbar=0} = \tau$ &
 Free BRST complex \cite{CDG20} \\
B (inversion) &
 GK involution $=$ identity &
 Self-mirror \\
C (complementarity) &
 Cross-polarization &
 scalar $\kappa$-surface complementarity \\
D (modular char.) &
 $\Theta_{\beta\gamma}\big|_{g=1,\,n=0}
 = \kappa \cdot \lambda_1$ &
 $\kappa(\cA_{\beta\gamma}) = 6\lambda^2 - 6\lambda + 1$;\; $=1$ at $\lambda = 0,1$;\; $= -\tfrac{1}{2}$ at $\lambda = \tfrac{1}{2}$ \\
H (chiral Hochschild) &
 Degree-preserving sub-MC &
 $\ChirHoch^n(\beta\gamma_\lambda)=0$ for
 $n\notin\{0,1,2\}$ \\
\bottomrule
\end{tabular}
\end{center}
The collision projection of the closed chiral algebra has
$r^{\mathrm{coll}}_{\beta\gamma}=0$, and its $bc$ Koszul partner has
$r^{\mathrm{coll}}_{bc}=0$ by the same simple-pole convention. The
nonzero rational expression \eqref{eq:betagamma-r-matrix} is instead
the open line evaluation $r^{\mathrm{line}}_{\beta\gamma}(z)$ of
$\tau_{\mathrm{line}}$; it is not a genus-$0$ closed collision
residue and it is not the closed chiral Koszul dual
$A^!_{\mathrm{ch}}\simeq bc_\lambda$. Because the local
$\beta\gamma$ OPE is free, closed higher local vertices vanish; the
standard-family quartic datum is a transferred modular contact
coefficient.
\end{remark}

\section{Synthesis}
\label{sec:bg-supplement}

\begin{remark}[$\beta\gamma$-system as K3 building block]
\label{rem:bg-1}
The $\beta\gamma$-ghost system developed here, the class-C witness in
the four-class G/L/C/M Dichotomy with shadow-tower depth $r_{\max} = 4$,
enters the K3 chiral bialgebra $\mathbf H_{\Delta_5}$ as a building block. The class-$\mathcal S$
$A_1$-on-$\Sigma_{0,24}$ VOA decomposes, in a distinguished chamber of
its moduli space, as a tensor product of $24$ pairs of $\beta\gamma$
ghost systems, one per puncture of $\Sigma_{0,24}$. Concretely: at each
Kodaira I$_1$ fibre $p_a$ the half-BPS line operator $L_a^{\mathrm{half\text{-}BPS}}$
(Vol.~II, \S\ref{V2-sec:lineop-ext-supplement}) is generated by a
$\beta\gamma$ pair with weights adapted to the Mukai vector at $p_a$,
and the Beem--Lemos--Liendo--Peelaers--Rastelli--van~Rees 2015 chiral
algebra is the $M_{24}$-invariant quotient
\[
\chi[\mathcal T_{A_1,\Sigma_{0,24}}] \;=\; \Bigl(\bigotimes_{a=1}^{24} (\beta\gamma)_a\Bigr)^{M_{24}}.
\]
The Frenkel--Kac 1980 vertex-operator closure of $\mathcal F_{K3}$ then
realises this $24$-fold $\beta\gamma$-product after the Mukai-lattice
vertex operators are inserted. Primary-lit: Friedan--Martinec--Shenker
1986, Frenkel--Kac 1980, Beem et al.\ 2015.
\end{remark}
